\documentclass{report}
\usepackage{blindtext} % for table of contents and other niceseties.
\usepackage{amsmath}
\usepackage{amsfonts} % things such as \mathbb{R}, or \Z. See https://en.wikipedia.org/wiki/List_of_mathematical_symbols_by_subject
\usepackage{amstext} % for \text macro.
\usepackage{array} % for tables.
\usepackage{float} % here for H placement parameter
\usepackage{hyperref} % for hyperlinks.
\usepackage{listings} % for code.
\usepackage[makeroom]{cancel} % crossing out terms in math https://tex.stackexchange.com/questions/75525/how-to-write-crossed-out-math-in-latex
\usepackage[shortlabels]{enumitem} % enumerate with letters.
\usepackage{tikz}
\usepackage{tikz-3dplot}
\usepackage{braket}
\usepackage{booktabs} % for \addlinespace macro


% References.
\usepackage[
backend=biber,
style=alphabetic,
sorting=ynt
]{biblatex}
\addbibresource{references.bib}


% Make greek letters bold too.
\usepackage{bm}
\newcommand{\vect}[1]{\boldsymbol{\mathbf{#1}}}

% The following is meant to help us get the cursive r script thats used in electrodynamocs (the separation vector).
% taken from https://www.reed.edu/physics/faculty/griffiths/
\usepackage{graphicx}
\def\rcurs{{\mbox{$\resizebox{.16in}{.08in}{\includegraphics{ScriptR}}$}}}
\def\brcurs{{\mbox{$\resizebox{.16in}{.08in}{\includegraphics{BoldR}}$}}}
\def\hrcurs{{\mbox{$\hat \brcurs$}}}

% \arg already exists.
\DeclareMathOperator\Arg{Arg}

% Create an alternate definition of \Re, because the gothic lettering looks weird.
\DeclareMathOperator\re{Re}
\DeclareMathOperator\im{Im}



% Make paragrphs show in the table of contents.
\setcounter{secnumdepth}{4}
\setcounter{tocdepth}{4}


\title{Notes on Electrodynamics}

\numberwithin{equation}{section}

% There is a \pmod for the "mod" symbol, but not one for the "div" symbol.
% Thus we are making it up.
\makeatletter
\newcommand*{\bdiv}{%
  \nonscript\mskip-\medmuskip\mkern5mu%
  \mathbin{\operator@font div}\penalty900\mkern5mu%
  \nonscript\mskip-\medmuskip
}
\makeatother

\begin{document}

\maketitle
\tableofcontents

\section{Vector Analysis}

\subsection{Vector Algebra}

Face diagonals

\tdplotsetmaincoords{70}{120}
\begin{tikzpicture}
		[scale=2,
            tdplot_main_coords,
			cube/.style={thick,black},
			axis/.style={->,black,thick}]

	%draw the axes
	\draw[axis] (0,0,0) -- (1.5, 0,   0) node[anchor=west]{$x$};
	\draw[axis] (0,0,0) -- (0,   1.5, 0) node[anchor=west]{$y$};
	\draw[axis] (0,0,0) -- (0,   0,   1.5) node[anchor=west]{$z$};

	%draw the top and bottom of the cube.
	\draw[cube] (0,0,0) -- (0,1,0) -- (1,1,0) -- (1,0,0) -- cycle;
	\draw[cube] (0,0,1) -- (0,1,1) -- (1,1,1) -- (1,0,1) -- cycle;
	
	%draw the edges of the cube
	\draw[cube] (0,0,0) -- (0,0,1);
	\draw[cube] (0,1,0) -- (0,1,1);
	\draw[cube] (1,0,0) -- (1,0,1);
	\draw[cube] (1,1,0) -- (1,1,1);

    \coordinate (O) at (0,0,0);

    % Vector A.
    \pgfmathsetmacro{\ax}{1}
    \pgfmathsetmacro{\ay}{0}
    \pgfmathsetmacro{\az}{1}
    \coordinate (A) at (\ax,\ay,\az);
    \draw[->, very thick] (O) -- (A) node [midway, below left] {$\mathbf{A}$};
    \draw plot [mark=*, mark size=0.2] (A) node [above, left] {\scriptsize$(1,0,1)$};

    % Vector B.
    \pgfmathsetmacro{\bx}{0}
    \pgfmathsetmacro{\by}{1}
    \pgfmathsetmacro{\bz}{1}
    \coordinate (B) at (\bx,\by,\bz);
    \draw[->, very thick] (O) -- (B) node [midway, below right] {$\mathbf{B}$};
    \draw plot [mark=*, mark size=0.2] (B) node [above, right] {\scriptsize$(0,1,1)$};

    % Angle.
    \tdplotdefinepoints(0,0,0)(\ax,\ay,\az)(\bx,\by,\bz)
    \tdplotdrawpolytopearc[-, very thick]{0.5}{above}{$\theta$}  
\end{tikzpicture}

Body diagonals

\tdplotsetmaincoords{70}{120}
\begin{tikzpicture}
		[scale=2,
            tdplot_main_coords,
			cube/.style={thick,black},
			axis/.style={->,black,thick}]

	%draw the axes
	\draw[axis] (0,0,0) -- (1.5, 0,   0) node[anchor=west]{$x$};
	\draw[axis] (0,0,0) -- (0,   1.5, 0) node[anchor=west]{$y$};
	\draw[axis] (0,0,0) -- (0,   0,   1.5) node[anchor=west]{$z$};

	%draw the top and bottom of the cube.
	\draw[cube] (0,0,0) -- (0,1,0) -- (1,1,0) -- (1,0,0) -- cycle;
	\draw[cube] (0,0,1) -- (0,1,1) -- (1,1,1) -- (1,0,1) -- cycle;
	
	%draw the edges of the cube
	\draw[cube] (0,0,0) -- (0,0,1);
	\draw[cube] (0,1,0) -- (0,1,1);
	\draw[cube] (1,0,0) -- (1,0,1);
	\draw[cube] (1,1,0) -- (1,1,1);

    \coordinate (O) at (0,0,0);

    % Vector A.
    \pgfmathsetmacro{\ax}{1}
    \pgfmathsetmacro{\ay}{1}
    \pgfmathsetmacro{\az}{1}
    \coordinate (A) at (\ax,\ay,\az);
    \draw[->, very thick] (O) -- (A) node [above] {$\mathbf{A}$};
    \draw plot [mark=*, mark size=0.2] (A) node [above, right] {\scriptsize$(1,0,1)$};

    % Vector B.
    \pgfmathsetmacro{\bx}{1}
    \pgfmathsetmacro{\by}{1}
    \pgfmathsetmacro{\bz}{0}
    \coordinate (B) at (\bx,\by,\bz);
    \draw[->, very thick] (0,0,1) -- (B) node [below] {$\mathbf{B}$};
    \draw plot [mark=*, mark size=0.2] (B) node [below, right] {\scriptsize$(0,1,1)$};

    % Angle.
    \tdplotdefinepoints(0.5,0.5,0.5)(\ax,\ay,\az)(\bx,\by,\bz)
    \tdplotdrawpolytopearc[-, very thick]{0.5}{right}{$\theta$}  
\end{tikzpicture}
\section{Potentials}

\subsection{Multipole Expansion}


$$
\brcurs = |\mathbf{r} - \mathbf{r^\prime} |
$$
Where $\brcurs$ denotes the vector from a source point $\mathbf{r^\prime}$ to a
field point $\mathbf{r}$.
\chapter{Introduction to Electrostatics}


%%%%%%%%%%%%%%%%%%%%%%%%%%%%%%%%%%%%%%%%%%%%%%%%%%%%%%%%%%%%%%%%%%%%%%%
\section{Charge distributions}


%%%%%%%%%%%%%%%%%%%%%%%%%%%%%%%%%%%%%%%%%%%%%%%%%%%%%%%%%%%%%%%%%%%%%%%
\subsection{Surface elements $da$}

The following logic comes from Griffith's electrodynamics chapter 1.4.

In curvilinear coordinates, an infinitesimal displacement in the $\vect{\hat{r}}$ direction is
simply $dr$.
$$
dl_r = dr
$$

An infinitesimal element of length in the $\vect{\hat{\theta}}$ is
$$
dl_\theta = r d\theta
$$

And an infinitesimal element of length along the $\vect{\hat{\phi}}$ direction
is
$$
dl_\phi = r \sin\theta d\phi
$$

$\theta$ is the polar angle, the angle from the $z$ axis, and its range is $[0, \pi]$.
$\phi$ is the azimuthal angle, which is the angle around from the $x$ axis and its range is $[0, 2\pi]$.

The general infinitesimal displacement becomes
$$
d\vect{l} =
dr \vect{\hat{r}}
+ r d\theta \vect{\hat{\theta}}
+ r \sin\theta d\phi \vect{\hat{\phi}}
$$

From there one can see how the infinitesimal volume element $d\tau$ is then
$$
d\tau = dl_r dl_\theta dl_\phi =
r^2 \sin\theta dr d\theta d\phi
$$

\textbf{However, the inifinitesimal area element depends on the orientation of the surface!}

For example, if we integrate over the surface of a sphere, where r stays contant but $\theta$
and $\phi$ vary,
$$
d\vect{a} = dl_\theta dl_\phi \vect{\hat{r}} = r^2 \sin\theta d\theta d\phi \vect{\hat{r}}
$$

Where as is the surface lies in the x-y plane, so that $\theta$ is contant, and let's say $\theta = \phi/2$,
while $r$ and $\phi$ vary,
$$
d\vect{a} = dl_r dl_\phi \vect{\hat{\theta}}
= r \sin\theta dr d\phi \vect{\hat{\theta}}
= r dr d\phi \vect{\hat{\theta}}
$$

\textbf{The area element on the sphere can be calculated from the cross products of other two elements,}

You should also, definetely, see
\href{https://en.wikipedia.org/wiki/Spherical_coordinate_system#Integration_and_differentiation_in_spherical_coordinates}{Integration and differentiation in spherical coordinates}.
\\~\\


\textbf{Example 1}

One good example to get practice is Griffith's Example 1.13: find the volume of a sphere of radius $R$.

As we saw above, the optimal way to an answer is to integrate the volume element $d\tau$ in spherical coordinates.

\begin{align*}
V = \int d\tau &=
\int_{r=0}^{R} \int_{\theta=0}^{\pi} \int_{\phi=0}^{2\pi} dr \, d\theta \, d\phi r^2 \sin\theta \\
&=  \left( \int_{r=0}^{R} dr \,r^2 \right)
    \left( \int_{\theta=0}^{\pi} d\theta \, \sin\theta \right)
    \left( \int_{\phi=0}^{2\pi} d\phi \right) \\
&= \left( \frac{1}{3} R^3 \right) \left( -\cos\theta \Big|_{0}^{\pi} \right) \left( 2\pi \right) \\
&= \frac{4}{3} \pi R^3
\end{align*}



%%%%%%%%%%%%%%%%%%%%%%%%%%%%%%%%%%%%%%%%%%%%%%%%%%%%%%%%%%%%%%%%%%%%%%%
\subsection{Fundamental Theorems} \label{section:funamental-theorems}

This presentation is a quick summary of Griffith's 1.3.2 to 1.3.5.

The fundamental theorem of calculus states

$$
\int_{a}^{b} dx \, \left(\frac{df}{dx}\right) = f(b) - f(a)
$$

The fundamental theorem for gradients is

$$
\int_{\vect{a}}^{\vect{b}} \left(\nabla T\right) \cdot d\vect{l}
= T(\vect{b}) - T(\vect{a})
$$

Here $T$ is a scalar-values function, $T : \mathbb{R}^3 \rightarrow \mathbb{R}$, and the application of the
$\nabla$ operator ($\nabla = \left< \partial_x, \partial_y, \partial_z \right>$) makes $\nabla T$ a vector-valued function,
$\nabla T : \mathbb{R}^3 \rightarrow \mathbb{R}^3$.
The gradient points in the direction of maximum increase of the function $T$ and it expreses the rate of change of $T$
along a given direction.

The interesting bit of our fundamental theorem for gradients is that $\int_{\vect{a}}^{\vect{b}} \left(\nabla T\right) \cdot d\vect{l}$
is path independent.
And by consequence, $\oint \left(\nabla T\right) \cdot d\vect{l} = 0$
(if you go up the stairs, measure how many meters above sea-level you are, then go back down, and come back up, then your
change is 0).
\\

The fundamental theorem for divergences (Gauss's of Green's theorem) states that

$$
\int_{\mathcal{V}} \left(\nabla\cdot \vect{v}\right) \, d\tau
= \oint_{\mathcal{S}} \vect{v}\cdot d\vect{a}
$$

For the two previous fundamental theorems we saw that the boundary of a line was two points, and now we see that the boundary
of a volume $\mathcal{V}$ is the closed region $\mathcal{S}$.
Accordingly, if we sum how much quantity $\vect{v}$ spreads out over an entire volume, this is equal to how much
is $\vect{v}$ is crossing through the closed enclosing surface.
\\

The fundamental theorem for curls (Stoke's theorem) states that

$$
\int_{\mathcal{S}} \left( \nabla\times\vect{v} \right) \cdot d\vect{a}
= \oint_{\mathcal{P}} \vect{v} \cdot d\vect{l}
$$

Now the boundary of a region $\mathcal{S}$ is the perimeter (a line) $\mathcal{P}$.
Here we see that the circulation of $\vect{v}$, $\oint_{\mathcal{P}} \vect{v} \cdot d\vect{l}$, over a perimeter,
is equal to the sum of the flux of the "twisting" of $\vect{v}$ over a surface.


%%%%%%%%%%%%%%%%%%%%%%%%%%%%%%%%%%%%%%%%%%%%%%%%%%%%%%%%%%%%%%%%%%%%%%%
\subsection{Work done by a particle}

\textbf{Taken from Kline, Chapter 7, Section 7, problem 3}

AM is a straight vertical line.
MO is a straight horizontal line. O is to the right of M.

A particle moving along AM is attracted to a fixed point O with a force that varies inversely
as the square of the distance from O.
Find the work done on the particle as it moves from A to a distance B (between A and M).

Suggestion: let $x$ be the distance from A to any point P on AB
and let $r$ be the variable distance from P to O.
The force $F$ attracting P to O is $k/r^2$.
But only the component of this force along AB serves to move the particle along AB.
This component is $(k/r^2) \cos OPM = (k/r^2) (c-x)/r$.
Then $dW/dx = k (c-x)/r^3$ and $r = \sqrt{p^2 + (c-x)^2}$.
\\

\begin{align*}
W &= \int_{A}^{B} \frac{dW}{dx} dx \\
&= \int_{A}^{B} \frac{k (c-x)}{r^3} \\
&= \int_{A}^{B} \frac{k (c-x)}{\left(p^2 + (c-x)^2\right)^{3/2}} dx \\
\end{align*}

The trick here is to note that if we chose $u = p^2 + (c-x)^2$, then $du = -2(c-x) dx$, so that
$$
\frac{k (c-x)}{\left(p^2 + (c-x)^2\right)^{3/2}} dx \rightarrow -\frac{k}{2}\frac{1}{u^{3/2}} du
$$

Since $\int u^{-3/2} du = -2 u^{-1/2} + C$.
So putting everything together,
\begin{align*}
W &= \int_{A}^{B} \frac{k (c-x)}{\left(p^2 + (c-x)^2\right)^{3/2}} dx \\
&= -\frac{k}{2} (-2) \frac{1}{u^{1/2}} \Big|_{A}^{B} \\
&= k \left( \frac{1}{\sqrt{p^2 + (c-B)^2}} - \frac{1}{\sqrt{p^2 + (c-A)^2}} \right) \\
&= \frac{k}{r_B} - \frac{k}{r_A}
\end{align*}

\textbf{Note:} the two interesting bits to note from this exercise are first, when we were looking for the component
of motion that contributed to the work being done, the vertix of the angle used was the point we were "on".
Second, note how our the u-substitution we used went, doing it the normal way may confuse you
and try to get you for somewhere where to use $u = p^2 + (c-x)^2$, which is not needed.
\\~\\


%%%%%%%%%%%%%%%%%%%%%%%%%%%%%%%%%%%%%%%%%%%%%%%%%%%%%%%%%%%%%%%%%%%%%%%
\section{Gauss's Law}



%%%%%%%%%%%%%%%%%%%%%%%%%%%%%%%%%%%%%%%%%%%%%%%%%%%%%%%%%%%%

\subsection{Solid Angle}

The following is all wikipedia.

Whereas an angle in radians, projected onto a circle, gives a length of a circular arc
on the circumference, a solid angle in steradians, projected onto a sphere,
gives the area of a spherical cap on the surface.

The point from which the object is viewed is called the apex of the solid angle,
and the object is said to subtend its solid angle at that point.

In SI units, a solid angle is expressed in a dimensionless unit called a \textbf{steradian}.

One steradian corresponds to one unit of area on the unit sphere surounding the apex,
so an object that blocks all rays from the apex would cover a number of steradians equal
to the total surface area of the unit sphere, $4\pi$.


Just like a planar angle in radians is the ratio of the length of an arc to its radius, $\theta = s/r$,
a solid angle in steradians is the ratio of the area covered on a sphere by an object to the area given
by the square of the radius of said sphere. The formula is
$$
\Omega = \frac{A}{r^2}
$$
where A is the spherical surface area and r is the radius of the considered sphere.

Any area on a sphere which is equal in area to the square of its radius,
when observed from its center, subtends precisely one steradian.

The solid angle of a sphere measured from any point in its interior is $4\pi$ sr (steridians),
and the solid angle subtended at the center of a cube by one of its faces is one-sixth of that,
or $2\pi/3$ sr.

In spherical coordinates there is a formula for the differential,
$$
d\Omega = \sin \theta d\theta d\varphi 
$$
where $\theta$ is the colatitude (angle from the North Pole - z-axis range $[0,\pi]$)
and $\varphi$ is the longitude (range $[0,2\pi]$).

The solid angle for an arbitrary oriented surface S subtended at a point P is equal to
the solid angle of the projection of the surface S to the unit sphere with center P,
which can be calculated as the surface integral:
$$
\Omega =\iint _{S} \frac{\hat{r}\cdot\hat{n}}{r^2} \,dS
= \iint_{S} \sin \theta \,d\theta \,d\varphi
$$

where $\hat{r} = \vec{r}/r$ is the unit vector corresponding to $\vec {r}$,
the position vector of an infinitesimal area of surface $dS$ with respect to point P,
and where $\hat{n}$ represents the unit normal vector to $dS$.
Even if the projection on the unit sphere to the surface S is not isomorphic,
the multiple folds are correctly considered according to the surface orientation
described by the sign of the scalar product $\hat{r}\cdot\hat{n}$.

Thus one can approximate the solid angle subtended by a small facet having flat surface area
$dS$, orientation $\hat{n}$, and distance r from the viewer as:
$$
d\Omega = 4\pi \left( \frac{dS}{A} \right) (\hat{r}\cdot\hat{n})
$$
where the surface area of a sphere is $A = 4\pi r^2$.


%%%%%%%%%%%%%%%%%%%%%%%%%%%%%%%%%%%%%%%%%%%%%%%%%%%%%%%%%%%%
\subsection{Explanation of Jackson}

If the electric field makes an angle $\theta$ with the unit normal, then the projection of the
infinitesimal area element along the normal of the surface is $\cos\theta da$.
Using the solid angle formula, we then get $d\Omega = \cos\theta da / r^2$ or
$r^2 d\Omega = \cos\theta da$.



%%%%%%%%%%%%%%%%%%%%%%%%%%%%%%%%%%%%%%%%%%%%%%%%%%%%%%%%%%%%%%%%%%%%%%%%%%%%%%%%%%%%%%%%%%%%%%%%%%%%%%%%%%%%%%%%%%%%%%
\section{Problems}

%%%%%%%%%%%%%%%%%%%%%%%%%%%%%%%%%%%%%%%%%%%%%%%%%%%%%%%%%%%%%%%
\subsection{Problem 1.2}
\label{jackson:problem-1.2}


The Dirac delta function in three dimensions can be taken as the improper limit as $\alpha \rightarrow 0$ of the
Gaussian function
$$
D\left(\alpha ; x, y, z \right) =
    \frac{1}{\left(\alpha \sqrt{2\pi}\right)^3}
    \exp{ \left[ -\frac{1}{2\alpha^2} \left( x^2 + y^2 + z^2 \right) \right] }
$$

Consider a general orthogonal coordinate system specified by the surfaces $u$ = constant, $v$ = constant,
$w$ = constant, with length elements $du/U$, $dv/V$, $dw/W$ in the three perpendicular directions.
Show that
$$
\delta\left( \mathbf{x} - \mathbf{x}^\prime \right) =
    \delta\left(u - u^\prime\right) \delta\left(v - v^\prime\right) \delta\left(w - w^\prime\right) \cdot UVW
$$

by considering the limit of the Gaussian above.
Note that as $\alpha \rightarrow 0$ only the infinitesimal length
element need be used for the distance between the points in the exponent.
\\

Browsing through the web, one will see the following defintions of the delta function,
\begin{align*}
\delta (x) &= \lim_{b\rightarrow 0} \frac{1}{|b|\sqrt{\pi}} \exp{ \left[ -\frac{1}{b^2} x^2 \right]} \\
    &= \lim_{b\rightarrow 0^+} \frac{1}{2\sqrt{\pi b}} \exp{ \left[ -\frac{1}{4b} x^2 \right]} \\
    &= \lim_{b\rightarrow 0} \frac{1}{b \sqrt{2\pi}} \exp{ \left[ -\frac{1}{2b^2} x^2 \right]}
\end{align*}

The thing to note here is how the arguments of the exponent make it into the constant that's infront of the exponent.
And the reason to even look at this now is because we found this amazingly interesting approach to this problem in
\href{https://www.wtamu.edu/~cbaird/courses.html}{West Texas CS Baird's} site:

We start with
$$
D\left(\alpha ; x, y, z \right) =
    \frac{1}{\left(\alpha \sqrt{2\pi}\right)^3}
    \exp{ \left[ -\frac{1}{2\alpha^2} \left( x^2 + y^2 + z^2 \right) \right] }
$$
and then we make a change of variables, a displacement if you will, by changing $x \rightarrow x - x^\prime$, etc.
This gives us
$$
D\left(\alpha ; x-x^\prime, y-y^\prime, z-z^\prime \right) =
    \frac{1}{\left(\alpha \sqrt{2\pi}\right)^3}
    \exp{ \left[ -\frac{1}{2\alpha^2} \left( (x-x^\prime)^2 + (y-y^\prime)^2 + (z-z^\prime)^2 \right) \right] }
$$

Now comes the limit taking.
First off is whether the limit is some actual number orr just infinity (we are skipping over the question of existence).
Since an exponent grows faster than a ever-increasing number to a power,
$\lim_{\alpha\rightarrow \infty} \alpha^3 \cdot e^{-\alpha} = 0$,
then we would expect our function $D$ to be bounded.
But, anyway, if we take the limit as $\alpha\rightarrow 0$, and if we look for a case where $D\not\rightarrow 0$,
then we would want $x-x^\prime$ to "balance out" $\alpha$ just enough so that we have some non-zero bounded value.

Let's assume the above is possible and rewirte $D$ as
$$
D\left(\alpha ; x-x^\prime, y-y^\prime, z-z^\prime \right) =
    \frac{1}{\left(\alpha \sqrt{2\pi}\right)^3}
    \exp{ \left[ -\frac{1}{2\alpha^2} \left( (dx)^2 + (dy)^2 + (dz)^2 \right) \right] }
$$

The infinitesimal displacements in the argument to the exponent look like an arc lenght, so let's rewrite it now as
$$
D\left(\alpha ; x-x^\prime, y-y^\prime, z-z^\prime \right) =
    \frac{1}{\left(\alpha \sqrt{2\pi}\right)^3}
    \exp{ \left[ -\frac{1}{2\alpha^2} \left( ds \right)^2 \right] }
$$

Since an arc length is a concept independent of the coordinate system, we can use our given general orthogonal coordinate system,
and rewrite our equation once more,
$$
D\left(\alpha ; u-u^\prime, v-v^\prime, w-w^\prime \right) =
    \frac{1}{\left(\alpha \sqrt{2\pi}\right)^3}
    \exp{ \left[ -\frac{1}{2\alpha^2} \left( 
        \left(\frac{du}{U}\right)^2 + \left(\frac{dv}{V}\right)^2 + \left(\frac{dw}{W}\right)^2 \right) 
    \right] }
$$

And if we reverse the process, by writing the differentials as differences,
$$
D\left(\alpha ; u-u^\prime, v-v^\prime, w-w^\prime \right) =
    \frac{1}{\left(\alpha \sqrt{2\pi}\right)^3}
    \exp{ \left[ -\frac{1}{2\alpha^2} \left(
        \left(\frac{u-u^\prime}{U}\right)^2 + \left(\frac{v-v^\prime}{V}\right)^2 + \left(\frac{w-w^\prime}{W}\right)^2 \right) 
    \right] }
$$

We can now look at the first couple definition of the delta function and make a final change of variahles.
For example, if we let $b_1 = \alpha U$, $b_2 = \alpha V$, and $b_3 = \alpha W$,
$$
D =
    \left[ \frac{
            \exp{ \left[ -\frac{
                \left(u-u^\prime\right)^2
            }{2b_{1}^{2}}  \right] }
        }{b_1 \sqrt{2\pi}}
        U
     \right]
     \left[ \frac{
            \exp{ \left[ -\frac{
                \left(v-v^\prime\right)^2
            }{2b_{2}^{2}}  \right] }
        }{b_2 \sqrt{2\pi}}
        V
     \right]
     \left[ \frac{
            \exp{ \left[ -\frac{
                \left(w-w^\prime\right)^2
            }{2b_{3}^{2}}  \right] }
        }{b_3 \sqrt{2\pi}}
        W
     \right]
$$

If we now take the limit of the $b_i$s, then we get our desired result.
\\~\\

Yet another way to look at this problem is as follows.
The general property of the delta function is that
$$
\int \int \int \delta(x) \delta(y) \delta(z) dx dy dz = 1
$$

If we use our limit,
$$
\int \int \int \lim_{\alpha\rightarrow 0} D(\alpha; x, y, z) dx dy dz = 1
$$

In our general orthogonal coordinate system the above would look as
$$
\int \int \int \lim_{\alpha\rightarrow 0} D(\alpha; x, y, z \rightarrow u, v, w) \frac{du}{U} \frac{dv}{V} \frac{dw}{W} = 1
$$

Note that we do not know how to transform $D$, so we could wrap up our ignorance by defining an intermediate variable as such
$$
F(u,v,w) = \lim_{\alpha\rightarrow 0} D(\alpha; x, y, z \rightarrow u, v, w) \frac{1}{UVW}
$$
In which case we have
$$
\int \int \int F(u,v,w) du dv dw = 1
$$

If we conveniently generalize the above by taking into account the possibility of a translation - because we are integrating over
the entire universe anyway - then we have,
$$
\int \int \int F(u-u^\prime, v-v^\prime, w-w^\prime) du dv dw = 1
$$

Which has the exact same form as the expression that we started with, except we are looking at some other set of integration variables.
So we could make the claim that
\begin{align*}
F(u-u^\prime, v-v^\prime, w-w^\prime) &= \delta(u-u^\prime) \delta(v-v^\prime) \delta(w-w^\prime) \\
&= \lim_{\alpha\rightarrow 0} D(\alpha; x, y, z \rightarrow u, v, w) \frac{1}{UVW}
\end{align*}


And the most interesting part of this problem is the application!

Le'ts think again about spherical coordinates!
The lenght elements along the three orthogonal coordinates are $dr$, $rd\theta$, and $r\sin\theta d\phi$.
So $U=1$, $V=1/r$, and $W=1/r\sin\theta$.
So a three-dimensional delta function is spherical coordinates becomes

$$
\delta\left(\mathbf{x}-\mathbf{x}^\prime\right) =
\delta(r-r^\prime) \delta(\theta-\theta^\prime) \delta(\phi-\phi^\prime) \frac{1}{r^2 \sin\theta}
$$

One interesting way of writing the above is by using $\delta(\cos\theta)$.
Since we have the composition of a function as an argument to our delta function, we can use the following formula
$$
\delta\left(f(x)\right) = \sum_i \frac{1}{\left|\frac{df}{dx} (x_i)\right|} \delta(x-x_i)
$$
where the sum extends over all simple roots.

So,
$$
\delta(\cos\theta - \cos\theta^\prime)
= \sum \frac{1}{\left|-\sin\theta^\prime \right|} \delta(\theta-\theta^\prime)
= \frac{\delta(\theta-\theta^\prime)}{\sin\theta^\prime}
$$

Thus,
$$
\delta\left(\mathbf{x}-\mathbf{x}^\prime\right) =
\delta(r-r^\prime) \delta(\cos\theta-\cos\theta^\prime) \delta(\phi-\phi^\prime) \frac{1}{r^2}
$$

Now, in cylindrical coordinates, the lenght elements are $dr$, $rd\theta$, and $dz$, so
$U=1$, $V=1/r$, and $W=1$.
And
$$
\delta\left(\mathbf{x}-\mathbf{x}^\prime\right) =
\delta(r-r^\prime) \delta(\theta-\theta^\prime) \delta(z-z^\prime) \frac{1}{r}
$$


%%%%%%%%%%%%%%%%%%%%%%%%%%%%%%%%%%%%%%%%%%%%%%%%%%%%%%%%%%%%%%%%%%%%%%%%%%%%%%%%
\subsection{Problem 1.3}

Using Dirac Delta functions in the appropiate coordinate, express the following charge distributions as three-dimensional
charge densities $\rho(\mathbf{x})$.

\begin{enumerate}[(a)]
\item In spherical coordinates, a charge $Q$ uniformly distributed over a spherical shell of radius $R$.

\item In cylindrical coordinates, a charge $\lambda$ per unit length uniformly distributed over a cylindrical
surface of radius $b$.

\item In cylindrical coordinates, a charge $Q$ spread uniformly over a flat circular disc of negligible thickness and radius $R$.

\item The same as part (c), but using spherical coordinates.
\end{enumerate}
\chapter{Intro to QFT}

\section{Inivitation}

\subsection{Polarization Vectors}

In the invitation, we are presented with the polarization vector $\epsilon^\mu = (0, 1, i, 0)$.

To understand this notation let's go through some examples taken from
\href{https://scholar.harvard.edu/files/schwartz/files/lecture14-polarization.pdf}{Lecture 14: Polarization}.

We say that a plane wave is \textbf{linearly polarized} if there is no phase difference between $E_x$ and $E_y$:
$$
\vect{E}_0 = \left( E_x, E_y, 0 \right)
$$

Then a linearly polarized plane wave in the x direction looks like $\vect{E} = E_0 e^{i(kz-wt)} (1, 0, 0)$.
And a linearly polarized plane wave in the y direction looks like $\vect{E} = E_0 e^{i(kz-wt)} (0, 1, 0)$.

The thing to remember is that we are implicitly only looking at the real part, so a linearly polarized wave in the x
direction is actually $(E_0 \cos(kz-wt), 0, 0)$.
\\

\textbf{Circular polarization} is when the electric field components are one-quarter out of phase ($\pi /2$).
Then the field can be written as,

\begin{align*}
\vect{E}_0 &= \left( E_0, E_0 e^{i\pi/2}, 0 \right) \\
&= \left( E_0, iE_0, 0 \right) \\
&= E_0 \left( e^{i(kz-wt)}, ie^{i(kz-wt)}, 0 \right)
\end{align*}

And since we only care about the real parts,
$$
\vect{E}_0 = \left( cos(kz-wt), -sin(kz-wt), 0 \right)
$$

This is interesting because if you jump to the wikipedia page for "List of trigonometric identities"
and look for the "Shift by one quarter period" table, you'll see that
$$
\sin(\theta + \frac{\pi}{2}) = \cos(\theta)
$$
and
$$
\cos(\theta + \frac{\pi}{2}) = -\cos(\theta)
$$

(A shit by a quarter wavelength is essentially differentiation!)

But here is the catch and the connection with P\&S: we can also get the same result if we have
$$
\vect{E} = E_0 \left( \cos(wt-kz)\hat{x} + \sin(wt-kz)\hat{y} \right)
$$
We can also write is as,
$$
E_0 e^{i(wt - kz)} (0, 1, i, 0)
$$
and remembering that you only care about the real part.
If you do so, you'll end up with terms such as $\cos \hat{x} - \sin \hat{y}$ which just so happen to again be the
a $\sin$ and a $\cos$ (our original plane wave components) shifted by $\pi/2$.




\subsection{Cross Sections}

The next thing we want to document is how to solve the differential cross section (per unit solid angle).
The expression given was

$$
\frac{d\sigma}{d\Omega} = \frac{\alpha^2}{4E_{cm}^{2}} (1 + \cos^2 \theta)
$$

That when integrated gives the total cross section
$$
\sigma_{total} = \frac{4\pi\alpha^2}{3E_{cm}^{2}}
$$

The trick here is to identity $d\Omega = \sin\theta \, d\theta \, d\phi$.
So essentially the problem is to integrate

\begin{align*}
\int d\Omega \, \left( 1 + \cos\theta \right) &=
    \int_{\phi=0}^{2\pi} \int_{\theta=0}^{\pi} \left( 1 + \cos\theta \right) \sin\theta \, d\theta \, d\phi \\
&= \left( \int_{\phi=0}^{2\pi} \int_{\theta=0}^{\pi} \sin\theta \, d\theta \, d\phi \right)
    + \left( \int_{\phi=0}^{2\pi} \int_{\theta=0}^{\pi} \cos\theta \sin\theta \, d\theta \, d\phi \right) \\
&= \left( 4\pi \right) + \left( -\int u^2 du \right) \\
&= \left( 4\pi \right) + \left( -\frac{1}{3} \cos^3 \theta \Big|_{0}^{\pi} du \right) \\
&= \left( 4\pi \right) \left( 2\pi \frac{8}{3} \right) \\
&= \frac{16\pi}{3}
\end{align*}

Note that we did a $u$-substitution in the second integral: $u = \cos\theta$, so $du = -\sin\theta \, d\theta$.







%%%%%%%%%%%%%%%%%%%%%%%%%%%%%%%%%%%%%%%%%%%%%%%%%%%%%%%%%%%%%%%%%%%%%%%%%%%%%%%%%%%%

\section{The Klein-Gordon Field}

\subsection{Klein-Gordon Inconsistensies}

The Schrodinger equation canr readily be abtain treating the energy and momentum as operators.
In quantum mechanics $E = i\partial_t$ and $p = -i\nabla$.
Using the realtionship $E = \frac{p^2}{2}$, we get


$$
i\partial_t \ket{\psi} = -\frac{1}{2m} \nabla^2 \ket{\psi}
$$

Note that we are still in natural units, otherwise there would be an $\hbar$.

The Klein-Gordon equation comes if we instead use $E^2 - p^2 = m^2$,

$$
-\partial_{t}^2 \ket{\psi} + \nabla^2 \ket{\psi} = m^2 \ket{\psi}
$$

And remember that $\partial^2 = \partial_{t}^{2} - \nabla^2$, so there comes $(\partial^2 - m)\psi = 0$ equation.



%%%%%%%%%%%%%%%%%%%%%%%%%%%%%%%%%%%%%%%%%%%%%%%%%%%%%%%%%%%%%%%%%%%%%%%%%%%%%%%
\subsubsection{Klein-Gordon Solution}

One method to solve the equation is to use Fourier transforms.
So we can begin by expressing the wave function $\phi(x,t)$ as a Fourier integral,

$$
\psi(x, t) =
\int \frac{d^4k}{\left(2\pi\right)^4} \tilde{\psi}(k, \omega) e^{-i(k\cdot x - \omega t)}
$$

Here, $\tilde{\psi}(k, \omega)$ is the Fourier transform of $\psi(x, t)$.
One thing that's often missed is the rgument of the exponent: if you look back to examples where we talk about plane waves
(just go back to the previous section), a wave with positive momentum moving in positive direction $x$ has the argument $wt - xk$.
And the convention is to hae such a wave in the Fourier transform.

With this Fourier transofrm in place, we take the Fourier transform of our entire differential equation to reduce
the solution into an algebraic problem.

For example, applying the Fourier transform to the time derivative goes as follows:

\begin{align*}
\mathcal{F}\{ \partial_t \psi \} &=
    \int \frac{d^4k}{(2\pi)^4} \partial_t \psi e^{-i(\vect{k}\cdot\vect{x} - \omega t)}
\rightarrow
\left[
    \begin{alignedat}{2}
        u  &= e^{-i(\vect{k}\cdot\vect{x} - \omega t)}                \quad & v  &= \psi \\
        du &= i\omega e^{-i(\vect{k}\cdot\vect{x} - \omega t)} \, dt  \quad & dv &= \partial_t \psi \, dt
    \end{alignedat}\,
\right] \\
&= \cancelto{0}{ \psi(\vect{x}, t) e^{-i(\vect{k}\cdot\vect{x} - \omega t)} } \Big|_{t=-\infty}^{t=\infty}
    - \int \frac{d^4k}{(2\pi)^4} \psi(\vect{x}, t) \left(i\omega\right) e^{-i(\vect{k}\cdot\vect{x} - \omega t)} \\
&= -i\omega \tilde{\psi} (\vect{k}, \omega)
\end{align*}

We made use of the integration by parts technique ($\int u \, dv = uv - \int v \, du$) to transfer the time derivative
on $\psi$ to the exponent term.

As per the boudnary term, we have a couple words to say: physically, the wave function (or even field) should vanish at infinity.
Mathematically, in order for the Fourier transform to converge to a value the function that is being "transformed", $\psi$,
must be a \textbf{rapidly decreasing function} (a function of \href{https://en.wikipedia.org/wiki/Schwartz_space}{Schwartz space}).

Similarly then,

\begin{align*}
\mathcal{F}\{ \partial_{t}^2 \psi \} &=
    \int \frac{d^4k}{(2\pi)^4} \partial_{t}^{2} \psi e^{-i(\vect{k}\cdot\vect{x} - \omega t)}
\rightarrow
\left[
    \begin{alignedat}{2}
        u  &= e^{-i(\vect{k}\cdot\vect{x} - \omega t)}                \quad & v  &= \partial_t \psi \\
        du &= i\omega e^{-i(\vect{k}\cdot\vect{x} - \omega t)} \, dt  \quad & dv &= \partial_{t}^{2} \psi \, dt
    \end{alignedat}\,
\right] \\
&= \cancelto{0}{ \partial _t \psi(\vect{x}, t) e^{-i(\vect{k}\cdot\vect{x} - \omega t)} } \Big|_{t=-\infty}^{t=\infty}
    - \int \frac{d^4k}{(2\pi)^4} \partial_t \psi(\vect{x}, t) \left(i\omega\right) e^{-i(\vect{k}\cdot\vect{x} - \omega t)} \\
&= -i\omega \mathcal{F}\{ \partial_t \psi \} \\
&= - \omega^2 \tilde{\psi} (\vect{k}, \omega)
\end{align*}

The boundary term again goes to zero here.
We will skip the physical argument and just mention that a rapidly decreasin function also requires all of its derivatives
tend to zero.



For the spatial derivatives, each $\partial_i$ will bring down a factor of $-i k_i$, resulting in a factor of
$- k_{i}^2$ for each $\partial_{i}^{2}$.
And this is something that always triped me up: one may be inclined to write $-k^2$ in the Fourier transform
but that same one ought to remember that $k$ is actually $\vec{k}$ (a vector) and that is why we must
write $-|k|^2$, because we want the [L-2] norm.

But anyways, the Fourier transform of the Laplacian is,

\begin{align*}
\mathcal{F}\{ \nabla \psi \} &=
    \int \frac{d^4k}{(2\pi)^4} \nabla \psi e^{-i(\vect{k}\cdot\vect{x} - \omega t)}
\rightarrow
\left[
    \begin{alignedat}{2}
        u  &= e^{-i(\vect{k}\cdot\vect{x} - \omega t)}                          \quad & v  &= \psi \\
        du &= -i\vect{k} e^{-i(\vect{k}\cdot\vect{x} - \omega t)} \, d\vect{x}  \quad & dv &= \nabla \psi \, d\vect{x}
    \end{alignedat}\,
\right] \\
&= \cancelto{0}{ \psi(\vect{x}, t) e^{-i(\vect{k}\cdot\vect{x} - \omega t)} } \Big|_{\vect{x}=-\infty}^{\vect{x}=\infty}
    - \int \frac{d^4k}{(2\pi)^4} \psi(\vect{x}, t) \left(-i\vect{k}\right) e^{-i(\vect{k}\cdot\vect{x} - \omega t)} \\
&= -i\vect{k} \tilde{\psi} (\vect{k}, \omega)
\end{align*}


and

\begin{align*}
\mathcal{F}\{ \nabla^2 \psi \} &=
    \int \frac{d^4k}{(2\pi)^4} \nabla^2 \psi e^{-i(\vect{k}\cdot\vect{x} - \omega t)}
\rightarrow
\left[
    \begin{alignedat}{2}
        u  &= e^{-i(\vect{k}\cdot\vect{x} - \omega t)}                          \quad & v  &= \nabla \psi \\
        du &= -i\vect{k} e^{-i(\vect{k}\cdot\vect{x} - \omega t)} \, d\vect{x}  \quad & dv &= \nabla^2 \psi \, d\vect{x}
    \end{alignedat}\,
\right] \\
&= \cancelto{0}{ \nabla\psi(\vect{x}, t) e^{-i(\vect{k}\cdot\vect{x} - \omega t)} } \Big|_{\vect{x}=-\infty}^{\vect{x}=\infty}
    - \int \frac{d^4k}{(2\pi)^4} \nabla\psi(\vect{x}, t) \left(-i\vect{k}\right) e^{-i(\vect{k}\cdot\vect{x} - \omega t)} \\
&= -i\vect{k} \mathcal{F}\{ \nabla \psi \} \\
&= - |\vect{k}|^2 \tilde{\psi} (\vect{k}, \omega)
\end{align*}



And since the above must hold throughout all of space, that's where
$$
\left(-\omega^2 + |\vect{k}|^2 + m^2 \right) \tilde{\psi}(\vect{k}, \omega) = 0
$$
comes from!

Now, in order for us to not have a trivial solution, $\psi(\vect{x}, t) = 0$, it must be so that
$\left(-\omega^2 + |\vect{k}|^2 + m^2 \right) = 0$, and so our dispersion relation comes about
$$
\omega^2 = |\vect{k}|^2 + m^2
$$
So $\omega = \pm \sqrt{ |\vect{k}|^2 + m^2 }$.
So essentially any function will do as long as the dispersion relation (momentum conservation) is respected.

A thing that people do, specially since we are talking about Fourier transforms is to define

$$
\psi(\vect{k}, \omega) = A(\vec{k}) \delta\left( \omega - \sqrt{ |\vect{k}|^2 + m^2 }\right) +
    B(\vec{k}) \delta\left( \omega + \sqrt{ |\vect{k}|^2 + m^2 }\right)
$$

Note that $A(\vec{k})$ and $B(\vec{k})$ are arbitrary functions that only depend on $\vect{k}$ since the
Dirac delta functions specifies the value for $\omega$.

With that, we can finally write a solution for the differential equation we started with,

\begin{align*}
\psi(\vect{x}, t) &=
    \int \frac{d^4k}{(2\pi)^4} \left[
        A(\vec{k}) \delta\left( \omega - \sqrt{ |\vect{k}|^2 + m^2 }\right) e^{-i(\vect{k}\cdot\vect{x} -\omega t)} +
        B(\vec{k}) \delta\left( \omega + \sqrt{ |\vect{k}|^2 + m^2 }\right) e^{-i(\vect{k}\cdot\vect{x} -\omega t)}
    \right] \\
&= \int \frac{d^4k}{(2\pi)^4} \left[
    A(\vec{k}) e^{-i(\vect{k}\cdot\vect{x} - \sqrt{ |\vect{k}|^2 + m^2 } t)} +
    B(\vec{k}) e^{-i(\vect{k}\cdot\vect{x} + \sqrt{ |\vect{k}|^2 + m^2 } t)}
\right]
\end{align*}

Hopefully thiis result makes sense: we obtained a family of function as solution and only initial conditions or
boundary conditions will result in a specific sort of function.




%%%%%%%%%%%%%%%%%%%%%%%%%%%%%%%%%%%%%%%%%%%%%%%%%%%%%%%%%%%%%%%%%%%%%%%%%%%%%%%
\subsubsection{Klein-Gordon Negative Density}


The expression to compute a probability current is derived by using the analogous continuity equation
from fluid dynamics

$$
\frac{\partial\rho}{\partial t} + \nabla\cdot \vect{j} = 0
$$


Using the Schrodinger equation and using the fact that $\rho = |\psi|^2 = \psi^* \psi$ one
can arrive at the probability density
$$
\rho = 
-\frac{i}{2m} \left( \psi^* \partial_t \psi - \psi \partial_t \psi^* \right)
$$

For the Klein Gordon equation, the way to massage it and get an expression for $\partial_t \rho$
and for its probability current is to:
take Klein-Gordon multiply it by $\psi^*$,
take the complex conjugate multiply it by $\psi$,
subtract the two and rearange to get something like the continuity equation.

Following those steps we have,

$$
\left( \partial_{t}^{2} - \nabla^2 + m^2 \right) \psi = 0
$$
and
$$
\left( \partial_{t}^{2} - \nabla^2 + m^2 \right) \psi^* = 0
$$

Multiplying them with $\psi^*$ and $\psi$ respectively we get,
$$
\psi^* \partial_{t}^{2} \psi - \psi^* \nabla^2 \psi + \psi^* m^2 \psi = 0
$$
and
$$
\psi \partial_{t}^{2} \psi^* - \psi \nabla^2 \psi^* + \psi m^2 \psi^* = 0
$$

Subtracting the latter from the former,
\begin{align*}
& \psi^* \partial_{t}^{2} \psi - \psi^* \nabla^2 \psi + \psi^* m^2 \psi -
    \psi \partial_{t}^{2} \psi^* + \psi \nabla^2 \psi^* - \psi m^2 \psi^* \\
&= \left( \psi^* \partial_{t}^{2} \psi - \psi \partial_{t}^{2} \psi^* \right)
    - \left( \psi^* \nabla^2 \psi - \psi \nabla^2 \psi^* \right)
    + \left( \psi^* m^2 \psi - \psi m^2 \psi^* \right) \\
&= \left( \psi^* \partial_{t}^{2} \psi - \psi \partial_{t}^{2} \psi^* \right)
    - \left( \psi^* \nabla^2 \psi - \psi \nabla^2 \psi^* \right) \\
&= \partial_t \left( \psi^* \partial_{t} \psi - \psi \partial_{t} \psi^* \right)
    -\nabla \cdot \left( \psi^* \nabla \psi - \psi \nabla \psi^* \right)
\end{align*}

The time derivative is then equated to $\partial_t \rho$ and the spatial derivatives to the current.
though one interesting tangent to take here is concerning the missing $\frac{i}{2m}$ factor that
these equations have.

If you remember, the imaginery part of of a complex number can be isolated by taking the different of it with its
complex conjugate: $\im{z} = \frac{1}{2i}(z - \overline{z})$. (Recall that if $z = a +ib$, then $\im{z} = b$, not $ib$.)
It just so happens that we are doing the same sort of operation here, so we can throw an $i/2$ factor into our equation
to ensure we get a real quantity.

For the $1/m$ factor we have to do some dimensional analysis.
To get the probability we have to integrate $\int d^3x \, |\psi|^2$, which is dimensionless.
So the $|\psi|^2$ term needs to cancel out the integration over space.
Hence $[|\psi|^2] = [L]^{-3} = [M]^{3}$.
From there we can say that $[\psi] = [L]^{-3/2}$ and
$[\psi^* \partial_{t}^2 \psi] = [L]^{-3/2}[L]^{-2}[L]^{-3/2} = [L]^{-5} = [M]^{5}$.

This same term we just evaluated needs to match the dimensions of
the time derivative, $[L]^{-1}$,a probability density, $[L]^{-3}$.
It should be then that $[\partial_t \rho] = [L]^{-4} = [M^4]$.
But hey, what a coincidence that if we were to divide $\psi^* \partial_{t}^2 \psi$ by a mass
that we would get just the right dimensions for our probability current!
And so it goes that the expression we are used to seeing turns out to be excused as such in the Klein-Gordon case.
\\


Suppose we had a $\psi =  e^{-i(\vect{k}\cdot\vect{x} -\omega t)}$ as solution.

Then,
\begin{align*}
\rho &=
    -\frac{i}{2m} \left( 
        e^{i(\vect{k}\cdot\vect{x} -\omega t)} \partial_t e^{-i(\vect{k}\cdot\vect{x} -\omega t)} -
        e^{-i(\vect{k}\cdot\vect{x} -\omega t)} \partial_t e^{i(\vect{k}\cdot\vect{x} -\omega t)} 
    \right) \\
&= -\frac{i}{2m} \left(
    e^{i(\vect{k}\cdot\vect{x} -\omega t)} \left(i\omega\right) e^{-i(\vect{k}\cdot\vect{x} -\omega t)} -
    e^{-i(\vect{k}\cdot\vect{x} -\omega t)} \left(-i\omega\right) e^{i(\vect{k}\cdot\vect{x} -\omega t)}
    \right) \\
&= -\frac{i}{2m} \left(i\omega\right) \left( 2 \right) \\
&= \frac{\omega}{m}
\end{align*}

And since $\omega = \pm \sqrt{|\vect{k}|^ + m^2}$, then the density can be negative!



%%%%%%%%%%%%%%%%%%%%%%%%%%%%%%%%%%%%%%%%%%%%%%%%%%%%%%%%%%%%%%%%%%%%%%%%%%%%%%%
\subsubsection{A Propagator}

Describes the amplitude for a particle to propagate from one point to another.

With a source $j(\vect{x}, t)$, the Klein-Gordon equation becomes
$$
\left( \partial^2 + m^2 \right) \psi(\vect{x}, t) = j(\vect{x}, t)
$$

In momentum space the equation is
$$
\left( -\omega^2 + |\vect{k}|^2 + m^2 \right) \tilde{\psi}(\vect{k}, \omega) = \tilde{j}(\vect{\vect{k}, \omega})
$$

$$
\tilde{\psi}(\vect{k}, \omega) = \frac{ \tilde{j}(\vect{\vect{k}, \omega}) }{ \omega^2 - |\vect{k}|^2 - m^2 + i\epsilon }
$$

$i\epsilon$ Feynman's prescription.

The propagator $D_F(x-y)$ is the inverse transform of
$\frac{ 1 }{ \omega^2 - |\vect{k}|^2 - m^2 + i\epsilon }$

$$
D_F(x-y) = \int \frac{d^4k}{(2\pi)^4}
    \frac{ e^{-i(\vect{k}\cdot\vect{x} -\omega t)} }{ \omega^2 - |\vect{k}|^2 - m^2 + i\epsilon }
$$






%%%%%%%%%%%%%%%%%%%%%%%%%%%%%%%%%%%%%%%%%%%%%%%%%%%%%%%%%%%%%%%%%%%%%%%%%%%%%%%%%%%%%%%%%%%
\subsection{Causality Arguments}

Consider the amplitude for a free particle to propagate from $\vect{x}_0$ to $\vect{x}$
$$
U(t) = \bra{\vect{x}} e^{-iHt} \ket{\vect{x}_0}
$$

In nonrelativistic quantum mechanics we have $E = \frac{1}{2m} p^2$, so
\begin{align}
U(t) &= \bra{\vect{x}} e^{-i(\vect{p}^2/2m)t} \ket{\vect{x}_0} \label{kg:nonrelativistic-01} \\
&= \int \frac{d^3p}{(2\pi)^3} \bra{\vect{x}} e^{-i(\vect{p}^2/2m)t} \ket{\vect{p}} \bra{\vect{p}}\ket{\vect{x}_0} \label{kg:nonrelativistic-02} \\
&= \int \frac{d^3p}{(2\pi)^3} e^{-i(\vect{p}^2/2m)t} e^{i\vect{p}\cdot(\vect{x}-\vect{x}_0)} \label{kg:nonrelativistic-03} \\
&= \left( \frac{m}{2\pi i t} \right)^{3/2} e^{im (\vect{x}-\vect{x}_0)^2 /2t} \label{kg:nonrelativistic-04}
\end{align}

The couple facts used in the above integral go as follows.

First, this comes from the section "outer products" in
\href{https://en.wikipedia.org/wiki/Bra%E2%80%93ket_notation}{Wikipedia: Bra-Ket Notation},
we have to keep in mind that $\ket{\psi}\bra{\psi}$ defines an \textbf{outer product}.
In a finite-dimensional vector space the outer product is defined as

$$
\ket{\phi}\bra{\psi}
=
\begin{pmatrix}
    \phi_1 \\
    \phi_2 \\
    \vdots \\
    \phi_N 
\end{pmatrix}
\begin{pmatrix}
    \psi_{1}^{*} & \psi_{2}^{*} & \ldots & \psi_{N}^{*}
\end{pmatrix}
=
\begin{pmatrix} 
    \phi_1 \psi_{1}^{*} & \phi_1 \psi_{2}^{*} & \dots  & \phi_1 \psi_{N}^{*} \\
    \phi_2 \psi_{1}^{*} & \phi_2 \psi_{2}^{*} & \dots  & \phi_2 \psi_{N}^{*} \\
    \vdots              & \vdots              & \ddots & \vdots              \\
    \phi_N \psi_{1}^{*} & \phi_N \psi_{2}^{*} & \dots  & \phi_N \psi_{N}^{*} \\
\end{pmatrix}
$$

One of the uses of the outer product is to construct \textbf{projection operators}. Given a ket $\ket{\psi}$
of norm 1, the orthogonal projection onto the subspace spanned by $\ket{\psi}$ is $\ket{\psi}\bra{\psi}$.


The "Unit operator" section on \href{https://en.wikipedia.org/wiki/Bra%E2%80%93ket_notation}{Wikipedia: Bra-Ket Notation},
also has this: if we have a complete orthonormal basis $\{ e_i | i\in\mathbb{N} \}$,
functional analysis tells us that any $\ket{\psi}$ can also be written as
$$
\ket{\psi} =
\sum_{i\in\mathbb{N}} \bra{e_i}\ket{\psi} \ket{e_i}
$$

This is how we "project" $\psi$ into a new basis.

It mentions that it can also be shown that
$$
\mathbb{I} = \sum_{i\in\mathbb{N}} \ket{e_i} \bra{\psi}
$$

There is also this result called
\href{https://en.wikipedia.org/wiki/Borel_functional_calculus#Resolution_of_the_identity}{resolution of the identity in Borel functional calculus}
that allows us to generalized this result to the continuous case,

$$
\mathbb{I} = \int dx \, \ket{x} \bra{x} = \int dp \, \ket{p} \bra{p}
$$

The difference with out notation is that we attach factor of $1/(2\pi)^d$ to our momentum-space integrals.
\\

The next bit of muscle memory to train is to remember that
$$
\delta^{4} (k) = \int \frac{d^4x}{(2\pi)^4} e^{ik\cdot x}
$$
and that
$$
\braket{x | x^\prime} = \delta(x - x^\prime)
$$

All this to say that to go from \ref{kg:nonrelativistic-01} to \ref{kg:nonrelativistic-02},
we chanegd the basis to the momentum space.
\section{Integrals in QFT}

\subsection{Kinetic operator and a Green Function} \label{integral:kinetic-operator}

This problem is part of
\href{https://arxiv.org/abs/physics/0609191}{MontePython: Implementing Quantum Monte Carlo using Python}.

There is a kinetic operator defined as
$$
G_K \left(\mathbf{R}^\prime, \mathbf{R}, t\right) =
\frac{1}{\left(2\pi\right)^{3N}} \int e^{-i\mathbf{k}\mathbf{R}^\prime} e^{-Dtk^2} e^{-i\mathbf{k}\mathbf{R}} d\mathbf{k}
$$

This integral then is integrated and the following Green function is obtained
$$
G_K \left(\mathbf{R}^\prime, \mathbf{R}, t\right) =
\frac{1}{\left(4\pi Dt\right)^{3N/2}} e^{-\left(\mathbf{R}-\mathbf{R}^\prime\right) / 4Dt}
$$

And as we saw in \hyperref[jackson:problem-1.2]{Jackson problem 1.2}, the Green function we just got is equivalent to
$\delta\left(\mathbf{R} - \mathbf{R}^\prime\right)$.

But coming back to the integral, the trick is to complete the square!
\begin{align*}
G_K \left(\vect{R}^\prime, \vect{R}, t\right) &=
\frac{1}{\left(2\pi\right)^{3N}} \int e^{-i\vect{k}\vect{R}^\prime} e^{-Dtk^2} e^{-i\vect{k}\vect{R}} d\vect{k} \\
&= \frac{1}{\left(2\pi\right)^{3N}} \int e^{-Dtk^2 + -i\vect{k}\cdot\left(\vect{R}-\vect{R}^\prime\right)} d\mathbf{k}
\end{align*}

This last integral has the form $\int e^{ax^2 + bx} dx$, which can be solved as follows,
\begin{align*}
\int e^{ax^2 + bx} dx &= \int \exp\left\{ a\left(x^2 + \frac{b}{a}x\right) \right\} dx \\
&= \int \exp\left\{ a\left(x^2 + \frac{b}{a}x + \left(\frac{b}{2a}\right)^2 - \left(\frac{b}{2a}\right)^2 \right) \right\} dx \\
&= \int \exp\left\{ a\left(x + \frac{b}{2a}\right)^2 - \frac{b^2}{4a} \right\} dx \\
&= e^{-b^2/4a} \int e^{ a\left(x + \frac{b}{2a}\right)^2 } dx
\end{align*}

Now, let's look at $\int e^{ a\left(x + b/2a\right)^2 } dx$.
First, we need to make a transformation.
Let's define $u = x + b/2a$, then $du = dx$, the measure remains invariant (so do the limits of integration), so
$\int e^{ a\left(x + b/2a\right)^2 } dx = \int e^{ au^2 } du$.
Now, if we define $a = -c$, then 

\begin{align*}
\int e^{ a\left(x + b/2a\right)^2 } dx &= \int e^{ au^2 } du \\
&= \int e^{ -cu^2 } du \\
&= \sqrt{ \frac{\pi}{c} } = \sqrt{ \frac{\pi}{-a} } 
\end{align*}

This whole thing works if $a < 0$ but it turns out that this is also valid if $\re{(a)} \leq 0$ but $a \neq 0$.

Putting everything back together,
$$
\int e^{ax^2 + bx} dx = \sqrt{ \frac{\pi}{-a} } e^{-b^2/4a}
$$

Looking back at our original problem, we can chose $a = -Dt$ and $b = i \vect{r}$, where $\vect{r} = \vect{R}-\vect{R}^\prime$.
and so
\begin{align*}
G_K \left(\vect{R}^\prime, \vect{R}, t\right) &=
\frac{1}{\left(2\pi\right)^{3N}} \int e^{-Dtk^2 + -i\vect{k}\cdot\vect{r}} d\mathbf{k} \\
&= \frac{1}{\left(2\pi\right)^{3N}} \int e^{-Dt \left( k^2 + \frac{i\vect{k}}{Dt}\cdot\vect{r} \right)} d\mathbf{k} \\
&= \frac{1}{\left(2\pi\right)^{3N}} \int e^{-Dt \left( k^2 + \frac{i\vect{r}}{Dt}\cdot\vect{k} + \left(\frac{i\vect{r}}{2Dt}\right)^2 - \left(\frac{i\vect{r}}{2Dt}\right)^2 \right)} d\mathbf{k} \\
&= \frac{1}{\left(2\pi\right)^{3N}} \int e^{-Dt \left( k + \frac{i\vect{r}}{2Dt}\right)^2 - \frac{ \vect{r}^2 }{ 4Dt } } d\mathbf{k} \\
&= \frac{1}{\left(2\pi\right)^{3N}} e^{-\vect{r}^2 /4Dt } \int e^{-Dt \left( k + \frac{i\vect{r}}{2Dt}\right)^2 } d\mathbf{k} \\
\end{align*}

Again, since $\int e^{-ax^2} dx = \sqrt{ \frac{\pi}{a} }$,
then
\begin{align*}
\int e^{-Dt \left( k + \frac{i\vect{r}}{2Dt}\right)^2 } dk &= \sqrt{ \frac{\pi}{Dt} }
\end{align*}

Thus,
\begin{align*}
G_K \left(\vect{R}^\prime, \vect{R}, t\right) &=
\frac{1}{\left(2\pi\right)^{3N}} e^{-\vect{r}^2 /4Dt } \int e^{-Dt \left( k + \frac{i\vect{r}}{2Dt}\right)^2 } d\mathbf{k} \\
&= \frac{1}{\left(2\pi\right)^{3N}} e^{-\vect{r}^2 /4Dt } \left( \sqrt{ \frac{\pi}{Dt} } \right)^{3N} \\
&= \frac{1}{\left(4\pi^2\right)^{3N/2}} \frac{\pi^{3N/2}}{(Dt)^{3N/2}} e^{-\vect{r}^2 /4Dt } \\
&= \frac{1}{\left(4\pi Dt\right)^{3N/2}} e^{-\left(\vect{R} - \vect{R}^\prime\right)^2 /4Dt }
\end{align*}



%%%%%%%%%%%%%%%%%%%%%%%%%%%%%%%%%%%%%%%%%%%%%%%%%%%%%%%%%%%%%%%%%%%%%%%%%%%%%%%%%%%%%%%%%%%%%%%%%%%%%%%%%

\subsection{Klein-gordon Non-relativistic Amplitude} \label{ps-kg:nonrelativistic-propagator}

The following integral comes from Peskin and schroeder Section 2.1
\begin{align}
& \int \frac{d^3p}{(2\pi)^3} \, e^{-i(\vect{p}^2/2m)t} e^{i\vect{p}\cdot(\vect{x}-\vect{x}_0)} \label{kg:nonrelativistic-03} \\
&= \left( \frac{m}{2\pi i t} \right)^{3/2} e^{im (\vect{x}-\vect{x}_0)^2 /2t} \label{kg:nonrelativistic-04}
\end{align}

This integral follows the same procedure as \ref{integral:kinetic-operator}.

Here we have $a = -\frac{it}{2m}$ and $b = i(\vect{x}-\vect{x_0})$.
Since $\re{a} = -t/2m \leq 0$ and $a\neq 0$ when $t\neq 0$ oru solution is valid for $t\neq 0$.
Then

$$
\sqrt{ \frac{\pi}{-a} } = \sqrt{ \frac{2\pi m}{it} }
$$

and
\begin{align*}
-\frac{b^2}{4a} &= \frac{2m \left(\vect{x}-\vect{x}_0\right)^2}{-4it} \\
&= \frac{im \left(\vect{x}-\vect{x}_0\right)^2}{2t}
\end{align*}

So the solution is
\begin{align*}
& \frac{1}{(2\pi)^3} \left( \frac{2\pi m}{it} \right)^{3/2} e^{\frac{im \left(\vect{x}-\vect{x}_0\right)^2}{2t}} \\
&= \left( \frac{m}{2\pi it} \right)^{3/2} e^{im \left(\vect{x}-\vect{x}_0\right)^2 /2t}
\end{align*}
%%%%%%%%%%%%%%%%%%%%%%%%%%%%%%%%%%%%%%%%%%%%%%%%%%%%%%%%%%%%%%%%%%%%%%%%%%%%%%%%%%%%%%%%%%%%%%%%%%%%%%%%%
\chapter{Tensors} \label{chapter:tensors}

Working with a spacetime metric signature $(+ - - -)$,
$$
\begin{pmatrix}
    1 & 0  & 0  & 0  \\
    0 & -1 & 0  & 0  \\
    0 & 0  & -1 & 0  \\
    0 & 0  & 0  & -1 \\
\end{pmatrix}
$$

$x^\mu = (x^0, \vect{x})$ and $x_\mu = g_{\mu\nu} x^\nu = (x^0, -\vect{x})$.
To see this explciitly, let's carry out the summations.

When $\mu = 0$, $x_0 = g_{00}x^0 + g_{01}x^1 + g_{02}x^2 + g_{03}x^3 = x^0 + 0 + 0 +0 = x^0$.
For $\mu = 1$, $x_1 = g_{10}x^0 + g_{11}x^1 + g_{12}x^2 + g_{13}x^3 = 0 + x^1 + 0 +0 = -x^0$.
Similarly, when $\mu = 2$, $x_2 = -x^2$, and when $\mu =3$, $x_3 = -x^3$.

The contravariant components of the vector $x^\mu$ are just its components.
The covariant compoents $x_\mu$, defined to be equal to $x_\mu = x^\nu g_{\nu\mu}$, are actually the components
of the associated dual vector.
Why? well, in order to obtain $x_\mu$ we had to use a "non-degenerate hermitian form"
(something that looks like an inner product; though it isn't since this metric is not positive definite).
This makes it a dual vector by definition.
\\

Also note that $g_{\mu\nu}$ forms a \textbf{bilinear form} since $g(x_1, x_2) = g(x_2, x_1)$
(linearity on both sides of the non-degenerate hermitian form).


%%%%%%%%%%%%%%%%%%%%%%%%%%%%%%%%%%%%%%%
\subsection{Inner Product}

Why does $p \cdot x = g_{\mu\nu} p^\mu x^\nu = g^{\mu\nu} p_\mu x_\nu = p_\mu x^\nu = p^\mu x_\nu = p^0 x^0 - \vect{p}\cdot\vect{x}$?

To start, we ought to know that there is the implicit defintion of
$x_\mu = x^\nu g_{\nu\mu}$.
So the covariant vector is defined through the application of a non-degenerate hermitian form
and it constitutes a dual vector.
From there the rest is linear algebra.

The metric tensor defines the geometry of spacetime, including the way distances and angles are measured.
The implied summation, einstein summation, effectively 'weights' the components according to the geometry of spacetime.

$g_{\mu\nu}$ lowers indices, while $g^{\mu\nu}$ raises them.
$\delta^{\mu}_{\nu}$ is a diagonal matrix with ones on the diagonal and zeros elsewhere.
It selects the $\mu$-th component when used in a summation, acting like an identity element.

$g^{\mu\alpha} g_{\alpha\nu} = \delta^{\mu}_{\nu}$

$$
g^{\mu\alpha} g_{\alpha\nu} = \sum_{\alpha = 0} g^{\mu\alpha} g_{\alpha\nu} =
g^{\mu 0} g_{0 \nu} + g^{\mu 1} g_{1\nu} + g^{\mu 2} g_{2\nu} + g^{\mu 3} g_{3\nu}
$$

This is essentially a matrix multiplication.
$$
\begin{pmatrix}
    1 & 0  & 0  & 0  \\
    0 & -1 & 0  & 0  \\
    0 & 0  & -1 & 0  \\
    0 & 0  & 0  & -1 \\
\end{pmatrix}
\cdot
\begin{pmatrix}
    1 & 0  & 0  & 0  \\
    0 & -1 & 0  & 0  \\
    0 & 0  & -1 & 0  \\
    0 & 0  & 0  & -1 \\
\end{pmatrix}
=
\begin{pmatrix}
    1 & 0 & 0 & 0 \\
    0 & 1 & 0 & 0 \\
    0 & 0 & 1 & 0 \\
    0 & 0 & 0 & 1 \\
\end{pmatrix}
$$

Going component by component, take the second row and second column, $\mu=1$ and $\nu=1$.
The element for that entry is given by
$$
g^{1\alpha}g_{\alpha 1} =
g^{1 0} g_{0 1} + g^{1 1} g_{1 1} + g^{1 2} g_{2 1} + g^{1 3} g_{3 1}
= 0 + (-1)(-1) + 0 + 0 = 1
$$


%%%%%%%%%%%%%%%%%%%%%%%%%%%%%%%%%%%%%%%
\subsection{Derivatives}

The examples here come from \cite{tensor-calc}.
The aim here is to see a few examples to get used to the quantities we are going to be seeing later on.
\\

Some useful identities before we carry on:
\begin{enumerate}
    \item
    $$
    \frac{\partial x_p}{\partial x_q} = \delta_{pq}
    $$
\end{enumerate}



\textbf{Example 1}

If $a_{ij}$ are contants, to compute the partial derivative
$$
\frac{\partial}{\partial x_k} a_{ij} x_i x_j
$$

We'll firs write out the sums explciitly.
\begin{align*}
a_{ij} x_i x_j &= \sum_{i,j} a_{ij} x_i x_j
\end{align*}

Someone once happened to noticed that the sum could be broken into parts depending on when $x_{\{i,j\}} = x_k$, where $x_k$
is the one we want to differentiate with respect to,
\begin{align*}
\sum_{i,j} a_{ij} x_i x_j &=
    \sum_{i \neq k, j \neq k} a_{ij} x_i x_j +
    \sum_{i, j \neq k} a_{ij} x_i x_j + 
    \sum_{i \neq k, j} a_{ij} x_i x_j +
    \sum_{i=k, j=k} a_{ij} x_i x_j \\
&= C + \sum_{j \neq k} \left( a_{kj} x_j \right) x_k +
    \sum_{i \neq k} \left( a_{ik} x_i \right) x_k +
    a_{kk} x_k x_k \\
    &= C + \sum_{j \neq k} \left( a_{kj} x_j \right) x_k +
    \sum_{i \neq k} \left( a_{ik} x_i \right) x_k +
    a_{kk} (x_k)^2 \\
\end{align*}

We wrote the first sum, where no $x$s are equal to $x_k$ as $C$, because as you may guess this is going to go away with
any derivative.

Back to our original task,
\begin{align*}
\frac{\partial}{\partial x_k} a_{ij} x_i x_j &=
    0 + \sum_{j \neq k} a_{kj} x_j +
    \sum_{i \neq k} a_{ik} x_i +
    2 a_{kk} x_k \\
&= \sum_{j} a_{kj} x_j + \sum_{i} a_{ik} x_i \\
&= a_{kj} x_j + a_{ik} x_i \\
&= a_{ki} x_i + a_{ik} x_i \\
&= (a_{ik} + a_{ki}) x_i
\end{align*}



\textbf{Example 2}

Another way to work out
$$
\frac{\partial}{\partial x_k} a_{ij} x_i x_j
$$
goes as follows.

\begin{align*}
\frac{\partial}{\partial x_k} a_{ij} x_i x_j &= a_{ij} \frac{\partial}{\partial x_k} x_i x_j \\
&= a_{ij} \left( \frac{\partial x_i}{\partial x_k} x_j + x_i \frac{\partial x_j}{\partial x_k} \right) \\
&= a_{ij} \left( \delta_{ik} x_j + x_i \delta_{jk} \right) \\
&= a_{kj} x_j + a_{ik} x_i \\
&= (a_{ik} + a_{ki}) x_i
\end{align*}


\textbf{Example 3}

If $a_{ij} = a_{ji}$ are constants, compute
$$
\frac{\partial^2}{\partial x_k \partial x_l} a_{ij} x_i x_j
$$

Let's start with our previous result,
\begin{align*}
\frac{\partial^2}{\partial x_k \partial x_l} a_{ij} x_i x_j &=
    \frac{\partial}{\partial x_k} \left[ \frac{\partial}{\partial x_l} a_{ij} x_i x_j \right] \\
&= \frac{\partial}{\partial x_k} (a_{il} + a_{li}) x_i \\
&= 2 a_{li} \frac{\partial x_i}{\partial x_k} \\
&= 2 a_{li} \delta_{ik} \\
&= 2 a_{lk}
\end{align*}



%%%%%%%%%%%%%%%%%%%%%%%%%%%%%%%%%%%%%%%%%%%%%%%%%%%%%%%%%%%%%%%%%%%%%%%%%%
\subsection{Changing Variables: the Jacobian} \label{subsection:jacobian}


Promise we are not doing the conventional thing where we go from scratch but there is one more pattern we have to covered
now in order for some other equations to make sense and this is the Jacobian.


There is a result from calculus that is presented a bit like this (the following presentation comes from
\cite{jacobian}):


If we define the Jacobian of a transformation $x = g(u,v)$, $y = h(u,v)$ as
\begin{align*}
\frac{ \partial (x,y) }{ \partial (u,v) } &= 
\begin{vmatrix}
\dfrac{ \partial x }{ \partial u } & \dfrac{ \partial x }{ \partial v } \\ \addlinespace
\dfrac{ \partial y }{ \partial u } & \dfrac{ \partial y }{ \partial v }
\end{vmatrix} \\
&= \frac{ \partial x }{ \partial u } \frac{ \partial y }{ \partial v } -
    \frac{ \partial x }{ \partial v }\frac{ \partial y }{ \partial u }
\end{align*}

Then if we want to integrate $f(x)$ over some surface $R$, this region will become some other surface $S$
following our transformation and the integral can be computed as follows:
$$
\int \int_{R} f(x,y) dA
= \int \int_{S} f\left( g(u,v), h(u,v) \right) \Biggl| \frac{ \partial (x,y) }{ \partial (u,v) } \Biggl| d\overline{A}
$$

For example, when we go from Euclidean to polar coordinates, the transformation we use is $x = r\cos\theta$
and $y = r\sin\theta$ and the Jacobian is $r$ so
$dA = \left| \frac{\partial (x, y)}{\partial (r, \theta)} \right| dx dy = r dr d\theta$:

\begin{align*}
\frac{\partial (x, y)}{\partial (r, \theta)} &= 
\begin{vmatrix}
    \dfrac{ \partial x }{ \partial r } & \dfrac{ \partial x }{ \partial \theta } \\ \addlinespace
    \dfrac{ \partial y }{ \partial r } & \dfrac{ \partial y }{ \partial \theta }
\end{vmatrix} \\
& = \begin{vmatrix}
    \cos\theta & -r\sin\theta \\ \addlinespace
    \sin\theta & r\cos\theta
\end{vmatrix} \\
&= r \cos^2 \theta + r \sin^2 \theta \\
&= r
\end{align*}



In 3-dimensions this extends to
$$
\int \int \int_{V} f(x,y,z) dV
= \int \int \int_{B} f\left( g(u,v,w), h(u,v,w), k(u,v,w) \right) \Biggl| \frac{ \partial (x,y,z) }{ \partial (u,v,w) } \Biggl| d\overline{V}
$$

For the case of the transformation of Euclidean coordinates to spherical coordinates, the transformation of variables is
$x = r \sin\phi \cos\theta $, $y = r \sin\phi \sin\theta$, and $z = r\cos\phi$.
In that case the Jacobian is $r^2 \sin\phi$ and thus
$dV = \left| \frac{\partial (x, y, z)}{\partial (r, \theta, \phi)} \right| dx dy dz = r^2 \sin\phi dr d\phi d\theta$.


Besides being very widely used, the reason we mentioned it is because this is also that is very well used in
\cite{gifted-qft} to introduce the notation we see so much in QFT.
Specifically, the Jacobian as we have been writing it, can also be seen in expressions such as

$$
\overline{a}^{\mu} = \left( \frac{\partial \overline{x}^\mu}{\partial x^\nu} \right) a^\nu
$$

$$
\frac{\partial \phi}{\partial \overline{x}^\mu} =
    \left( \frac{\partial x^\nu}{\partial \overline{x}^\mu} \right) \frac{\partial \phi}{\partial x^\nu}
$$

The determinant of the Jacobian matrix is used specifically when you are dealing with volume transformations
(like integrating over a volume or changing the measure of integration) because it scales the volume elements according
to how the coordinate transformation stretches or compresses the space.

When transforming vectors, however, you're not scaling a volume but rather reorienting or rescaling each component
of the vector according to how the coordinate axes themselves have changed.
This is why the transformation involves a sum over the product of vector components and the appropriate Jacobian elements:

\begin{itemize}
\item Contravariant vectors (standard vectors): The components of these vectors are transformed by multiplying with the
    Jacobian matrix $\partial \overline{x}^\mu / \partial x^\nu$. Here, each new vector component is a linear combination
    of old components weighted by how much each new coordinate axis changes with respect to each old axis.
    \begin{itemize}
        \item These objects, such as displacement vectors, inherently "follow" the coordinate axes. As the axes stretch
        or rotate, so do the components of the vector. 
    \end{itemize}
\item Covariant vectors (gradients and one-forms): These vectors transform using the inverse of the Jacobian matrix
    $\partial x^\nu / \partial \overline{x}^\mu$. This reflects how changes in the old coordinates map to the new
    coordinates, suitable for objects that naturally pair with vectors, like gradients.
    \begin{itemize}
        \item They measure rates of change along these coordinates.
        \item A larger coordinate stretch means a smaller gradient in that direction to maintain the same rate of change.
    \end{itemize}
\end{itemize}




%%%%%%%%%%%%%%%%%%%%%%%%%%%%%%%%%%%%%%%%%%%%%%%%%%%%%%%
\subsection{From Tensors to Matrices}

Tensors can be expressed as matrices only once a basis has been chosen.
the example being given is with the angular momentum operator $L_z$ on the spherical harmonics $H_l (\mathbb{R}^3)$.

It says that $H_l (\mathbb{R}^3)$ is the set of all linear functions on $\mathbb{R}^3$ that

$$
\{ r Y_{m}^{l} \}_{ -1 \leq m \leq 1 } =
\Big\{ \frac{1}{\sqrt{2}} (x +iy), \, z, \, - \frac{1}{\sqrt{2}} (x -iy) \Big\}
$$

and $\{ x, y, z \}$ are both basis for the space.

Then having the angular momentum operator $L_z = -i ( x\partial_y - y \partial_x )$ on this space we got
$$
\frac{1}{\sqrt{2}} L_z (x + iy) = \frac{1}{\sqrt{2}} (x + iy)
$$
$$
L_z (z) = 0
$$
$$
-\frac{1}{\sqrt{2}} L_z (x -iy) = \frac{1}{\sqrt{2}} (x -iy)
$$

which implies that in the spherical harmonics basis

$$
L_z = 
\begin{pmatrix}
1 & 0 & 0 \\
0 & 0 & 0 \\
0 & 0 & -1
\end{pmatrix}
$$

That tensor representation of $L_z$ is implies by saying that
$T(e_i) = \sum^{n}_{j=1} T_{i}^{j} e_j$.


$$
\frac{1}{\sqrt{2}} L_z (x + iy)
\rightarrow
L_z \left( \frac{1}{\sqrt{2}} (x + iy) \right)
$$

$$
L_z (z) = 0
$$

$$
-\frac{1}{\sqrt{2}} L_z (x -iy)
\rightarrow
L_z \left( -\frac{1}{\sqrt{2}} (x + iy) \right)
$$

Since

$$
\{ r Y_{m}^{l} \}_{ -1 \leq m \leq 1 } =
\Big\{ \frac{1}{\sqrt{2}} (x +iy), \, z, \, - \frac{1}{\sqrt{2}} (x -iy) \Big\}
$$

then we can say that $e_1 = \frac{1}{\sqrt{2}} (x +iy)$, $e_2 = z$, and $e_3 = -\frac{1}{\sqrt{2}} (x -iy)$.

Meaning that $L_z (e_1) = e_1 = 1\cdot e_1 + 0\cdot e_2 + 0\cdot e_3$,
$L_z (e_2) = 0 = 0\cdot e_1 + 0\cdot e_2 + 0\cdot e_3$, and
$L_z (e_3) = -e_3 = 0\cdot e_1 + 0\cdot e_2 - 1\cdot e_3$.
From here, we can then adopt write this as a matrix

$$
T_{i}^{j}
=
L_z
=
\begin{pmatrix}
1 & 0 & 0 \\
0 & 0 & 0 \\
0 & 0 & -1
\end{pmatrix}
$$

The columns of a matrix representation of a linear operator correspond to the action of the operator on the basis vectors.
However, due to the specific diagonal form of the matrix in this example, both the row and column perspectives
give the same result, which can indeed be confusing.

The first column represents $L_z(e_1)$,
the second column represents $L_z(e_2)$,
the third column represents $L_z(e_3)$.

In physics speak this would be read as the wave functions
$\frac{1}{\sqrt{2}} (x +iy)$, $z$, and $- \frac{1}{\sqrt{2}} (x -iy)$
have eigenvalues 1, 0, and -1, respectively.



%%%%%%%%%%%%%%%%%%%%%%%%%%%%%%%%%%%%%%%%%%%%%%%%%%%%%%%%%%%%%%%%%%%%%%%%%%%%%%%%%%%%%%
\section{Transformation Matrices}

The previous section gave us a good place to begin talking about tensors.
Now let's actually talk about them.

Tensors have been historically defined based on how they transform.
Think of the conventional axis rotation formula in 2D.
If we rotate the axis by an angle $\theta$, then

$$
x^\prime = x \cos\theta - y \sin\theta
$$
and
$$
y^\prime = x \sin\theta + y \cos\theta
$$

As we saw in the previous section,
where we derived the matrix representation of the angular momentum operator,
the columns of a matrix representation of a multilinear map (tensor) correspond to the action of
the map (tensor) on the basis vectors.
Which is most commonly seen as

$$
\begin{pmatrix}
x^\prime \\
y^\prime
\end{pmatrix}
=
\begin{pmatrix}
\cos\theta & - \sin\theta \\
\sin\theta & \cos\theta
\end{pmatrix}
\begin{pmatrix}
x \\
y
\end{pmatrix}
$$

If we envision $x$ and $y$, and $x^\prime$ and $y^\prime$, and some basis sets,
such that $x \rightarrow e_1$ and $y \rightarrow e_2$ (similarly for the primed components).
We can rewrite the above as
$$
e_{1^\prime} = e_1 \cos\theta - e_2 \sin\theta
$$
and
$$
e_{2^\prime} = e_1 \sin\theta + e_2 \cos\theta
$$

Meaning that our problem becomes,
$$
\begin{pmatrix}
e_{1^\prime} \\
e_{2^\prime}
\end{pmatrix}
=
\begin{pmatrix}
\cos\theta & - \sin\theta \\
\sin\theta & \cos\theta
\end{pmatrix}
\begin{pmatrix}
e_1 \\
e_2
\end{pmatrix}
$$

Following the Einstein summation convention, we could write

$$
e_{i^\prime} = A_{i^\prime}^{j} e_j
$$

which also means that the original basis can be obtained as follows
$$
e_{i} = A_{i}^{j^\prime} e_{j^\prime}
$$

It is interesting to now quickly pursue the exercise of expressing our basis in terms of itself
as it will show us some properties about these numbers $A_{i^\prime}^{j}$ and $A_{i}^{j^\prime}$.


Combining the previous two expressions we get,
$$
e_{i} = A_{i}^{j^\prime} e_{j^\prime}
= A_{i}^{j^\prime} A_{j^\prime}^{k} e_k
$$

From here we can see that
$$
A_{i}^{j^\prime} A_{j^\prime}^{k} = \delta_{i}^{k}
$$

If we consider the same example but with the primed-basis.

$$
e_{i^\prime} = A_{i^\prime}^{j} e_j 
= A_{i^\prime}^{j} A_{j}^{k^\prime} e_{k^\prime}
$$

Meaning that

$$
A_{i^\prime}^{j} A_{j}^{k^\prime} = \delta_{i^\prime}^{k^\prime}
$$

So the different $A$s are inverse of one another.

It is also worth noting that these $A$s are tensor components.


In expressions such as $A_{i^{\prime}}^{j} A_{j}^{k^{\prime}}$, the indices $i^{\prime}$ and $k^{\prime}$ refer
to the primed basis, while the index $j$ refers to the unprimed basis.
This mixing of basis sets in a single expression is why the numbers $A_{i^{\prime}}^{j}$ and $A_{j}^{k^{\prime}}$
are not considered tensor components.

The transformation matrices $A_{i^{\prime}}^{j}$ and $A_{j}^{k^{\prime}}$ are not tensor components themselves,
but rather express the relationship between components in different bases.
The fact that these matrices are inverses of each other, which is a consequence of the transformation law for
tensor components.

Another useful thing to note is that the transformation for dual vectors looks the same but with uppoer indices.

$$
e^{i^\prime} = A_{j}^{i^\prime} e^j
$$
and
$$
e^{i} = A_{j^\prime}^{i} e^{j^\prime}
$$

With this info at hand, we can guesstimate the matrix transformation laws for tensors,

$$
A =
\begin{pmatrix}
A_{1}^{1^\prime} & A_{2}^{1^\prime} & \dots  & A_{n}^{1^\prime} \\
A_{1}^{2^\prime} & A_{2}^{2^\prime} & \dots  & A_{n}^{2^\prime} \\
\vdots           & \vdots           & \ddots & \vdots           \\
A_{1}^{n^\prime} & A_{2}^{n^\prime} & \dots  & A_{n}^{n^\prime}
\end{pmatrix}
$$

Its inverse is
$$
A^{-1}
=
\begin{pmatrix}
A_{1^\prime}^{1} & A_{2^\prime}^{1} & \dots  & A_{n^\prime}^{1} \\
A_{1^\prime}^{2} & A_{2^\prime}^{2} & \dots  & A_{n^\prime}^{2} \\
\vdots           & \vdots           & \ddots & \vdots           \\
A_{1^\prime}^{n} & A_{2^\prime}^{n} & \dots  & A_{n^\prime}^{n} \\
\end{pmatrix}
$$

These are unitary matrices
$$
A A^{-1} = A^{-1} A = I
$$

Equating the above with our 2D rotation example,
$$
A^{-1}
=
\begin{pmatrix}
\cos\theta & - \sin\theta \\
\sin\theta & \cos\theta
\end{pmatrix}
$$


%%%%%%%%%%%%%%%%%%%%%%%%%%%%%%%%%%%%%%%%%%%%%%%%%%%%%%%%%%%%%%%%%%%%%%%%%%%%%%%%%%%%%%
\section{Contravariant and Covariant Vectors}

One fun fact about vectors is that they can be seen as tensors of type (0,1), that is, they feed
on 1 dual vector.
A vector can be viewed as a tensor of type (0, 1) because it transforms according to the contravariant
transformation law under a change of coordinates. The key idea is that a vector represents a directed
quantity that "eats" or operates on dual vectors (covectors or linear functionals) to produce scalars.

This is mathematically said as,
$$
v^{i^\prime} = A^{i^\prime}_{j} v^j
$$
As we saw above, this is the transformation for dual vectors - expressions with contravariant indices.

And as we saw above, these transform according to the Jacobian.

A dual vector in turn transforms as
$$
f_{i^\prime} = A_{i^\prime}^{j} v_j
$$

Note how the transformation are the contrary to what we saw for the basis vectors!
Where as the components of the dual vector transform in the same way as the basis.

In matrix notation we have contravariant vectors transform as
$$
v_{\mathcal{B}^\prime} = A v_{\mathcal{B}}
$$

and covariant ones as
$$
f_{\mathcal{B}^\prime} = \left( A^{-1} \right)^T f_{\mathcal{B}}
$$

But remember that basis and dual basis transform in the opposite ways!
\\


Now, in quantum mechanics, if we consider an orthonormal basis ${\ket{e_{i^{\prime}}}}$ in the same Hilbert space.
And assume that both ${\ket{e_i}}$ and ${\ket{e_{i^{\prime}}}}$ are complete bases,
so there must be a unitary transformation $U$ that relates them.
We would write, $\ket{e_{i^{\prime}}} = U \ket{e_i}$ for this.
Consecuently, the dual basis transforms as follows, $\bra{ e_{i^\prime} } = \bra{e_i} U^\dagger $.

Here, $U$ plays the role of $A^{-1}$, or $A_{i^{\prime}}^{j}$ in the tensor formalism.

Following the previous approach, we get the following results for transformations into a primed basis set,
$$
\ket{e_{i^{\prime}}} = U \ket{e_i}
$$
and
$$
\bra{ e_{i^\prime} } = \bra{e_i} U^\dagger
$$

The matrix elements of $U$ in the $\ket{e_i}$ basis are given by:
$U_{i^{\prime}}^{j} = \bra{e_{i^{\prime}}} U \ket{e_j}$.
These matrix elements play the same role as the $A_{i^{\prime}}^{j}$ in your classical expression.






%%%%%%%%%%%%%%%%%%%%%%%%%%%%%%%%%%%%%%%%%%%%%%%%%%%%%%%%%%%%%%%%%%%%%%%%%%%%%%%%%%%%%%
\subsection{Pauli Matrices}

The Pauli matrices actually form a basis for $2 \times 2$ Hermitian matrices.
Let's see how: a generic $2 \times 2$ Hermitian matrix would look like

$$
\begin{pmatrix}
a & b\\
c & d
\end{pmatrix}
=
\begin{pmatrix}
a^* & c^* \\
b^* & d^*
\end{pmatrix}
$$

Meaning that $a, b \in \mathbb{R}$ and that $b = c^*$.
So we could write our Hermitian matrix as

$$
\begin{pmatrix}
a      & x - iy \\
x + iy & d
\end{pmatrix}
$$

Another requirement that is imposed on $H_2 (\mathbb{C})$ is that they must be unitary, that is $U^\dagger U = I$, so

\begin{align*}
\begin{pmatrix}
a      & x - iy \\
x + iy & d
\end{pmatrix}
\begin{pmatrix}
a      & x - iy \\
x + iy & d
\end{pmatrix}
&=
\begin{pmatrix}
a^2 + x^2 - y^2         & a(x - iy) + b(x - iy) \\
a(x + iy) + b(x + iy) & d^2 + x^2 - y^2
\end{pmatrix}
\\
&=
\begin{pmatrix}
1 & 0 \\
0 & 1
\end{pmatrix}
\end{align*}

One way to make this happen is to set $a = -d$, this will also have the happy side effect of meeting another requirement
for the Pauli matrices to be traceless.
So now we have our candidate bases that ought to look like

$$
\begin{pmatrix}
a      & x - iy \\
x + iy & -a
\end{pmatrix}
$$

This then means that our requirement for our matrix $H_2 (\mathbb{C})$ to be unitary simplifies to
$$
\begin{pmatrix}
a^2 + x^2 - y^2 & 0                \\
0               & a^2 + x^2 - y^2
\end{pmatrix}
=
\begin{pmatrix}
1 & 0 \\
0 & 1
\end{pmatrix}
$$

Here is where our 3 Pauli matrices comes!

$$
\sigma_1 = \sigma_x =  
\begin{pmatrix}
1 & 0 \\
0 & 1
\end{pmatrix}
$$
comes from assuming that $x=1$, so $a^2 + x^2 - y^2 = 1$ can be met by assuming that $a = y = 0$.

$$
\sigma_2 = \sigma_y =  
\begin{pmatrix}
0 & -i \\
i & 0
\end{pmatrix}
$$

comes from assuming that $y=1$, so $a^2 + x^2 - y^2 = 1$ can be met by assuming that $a = x = 0$.

Finally,
$$
\sigma_3 = \sigma_z =  
\begin{pmatrix}
1 & 0 \\
0 & -1
\end{pmatrix}
$$

Where we took $a=1$ and $x = y = 0$.




%%%%%%%%%%%%%%%%%%%%%%%%%%%%%%%%%%%%%%%%%%%%%%%%%%%%%%%%%%%%%%%%%%%%%%%%%%%%%%%%%%%%%
\section{Tensor Product}


By definition, a \textbf{vector} can be writen as
$$
v = v^i e_i
$$
See section 2.3.
$v^i$ are called the \textbf{contravariant} components of $v$.

From section 2.5, we see that we can also define the dual vector $e^i$ as
$$
e^i(v) = v^i
$$

A \textbf{dual vector} $f$ can then be defined as (remember that a dual vector is a linear functional
that takes a vector $v$ and returns a scalar)
$$
f(v)
= f(v^i e_i)
= v^i f(e_i)
= v^i f_i
$$
So while we refer to the compoents of a vector as $v^i$ on a basis $e_i$,
we refer to the components of a dual vector as $f_i$.
And you might have guessed that the dual vector itself can be writen as
$$
f_i = f_i e^i
$$

$f_i$ are the \textbf{covariant} components of $f$.

Given that a tensor is a multilinear map, the \textbf{linear operator}
can be seen as a case of a (1, 1) type tensor.
So for example, a linear operator $T$
$$
T(v)
= T(v^i e_i)
= v^i T(e_i)
= v^i T_{i}^{j} e_j
$$
Remember that we are using the Einstein summation convention so there is actually a double sum in the last
expression, $\sum_{i=1}^{n} \sum_{j=1}^{n} = \sum_{i,j}$.
Note that we also used the definition of linearity to "take out" the $v^i$ factor from the action of the
tensor.

Components of a tensor are the values of the tensor on a set of basis vectors.
For example, the components
$$
H_{i}^{j} = H(e_i, e^j)
$$

A vector can be considered a $(0, 1)$ type tensor (it feeds on a dual vector)
and its components transform like
$$
v^i = A_{j}^{i} v^j
$$
And a dual vector can be considered a $(1, 0)$ type tensor (it feeds on a vector)
and it compoents transform like
$$
f_i = A_{i}^{j} f_j
$$


%%%%%%%%%%%%%%%%%%%%%%%%%%%%%%%%%
\subsection{Formal definition}

Given two vectors $v \in V$ and $w \in W$,
the tensor product $v \otimes w$ is defined as the element of $V \otimes W$ such that
$$
(v \otimes w)(h, g) = h(v) g(w)
$$
For all $h \in V^*$ and $g \in W^*$.

If we consider some $T \in V \otimes W$ such that
$$
T(h, g) = h_i g_j T(e^i, f^j) = h_i g_j T^{i j}
$$


We defined $T^{i j}$ as $T^{i j} = T(e^i, f^j)$, $T$ which is a (0,2) tensor,
meaning that it's a multilinear map that takes two dual vectors as inputs and returns a scalar.
The vectors it takes as inputs are from the dual spaces (or covectors) $V^*$ and $W^*$.
$\{e^i\}$ is a basis for $V^*$ and $\{f^j\}$ is a basis for $W^*$.
Both sets of dual basis vectors, which is how we know $T^{ij}$ is a type (0,2) tensor.



The next term in our expression is $e_i \otimes f_j$.
This tensor product takes two dual vectors as inputs and returns a scalar,
$$
(e_i \otimes f_j)(v^*, w^*) = v^*(e_i) w^*(f_j)
$$

The action of the tensor product on these dual vectors is given by the action of each dual vector on its
corresponding basis vector.

This is exactly the behavior of a (0,2) tensor - it takes two dual vectors as inputs and returns a scalar.
$v^*(e_i)$ is a scalar because - a linear functional acting on a vector. Similarly,
$w^*(f_j)$.
\\

If we expand the dual vectors $v^*$ and $w^*$ in terms of their respective dual bases:
$$
v^* = v_i e^i
$$
$$
w^* = w_j f^j
$$

Then,

\begin{align*}
(e_i \otimes f_j)(v^*, w^*) &= v^*(e_i) w^*(f_j) \\
&= v_k e^k(e_i) \, w_l f^l(f_j) \\
&= v_i w_j
\end{align*}

This follows the same approach as in the book when deriving the expression after 3.39.
\\






Now let's put this into practice.


Suppose we have two vectors, $\vec{v} = (v_1, v_2)$ and $\vec{w} = (w_1, w_2, w_3)$.
The tensor product of these vectors, denoted $\vec{v} \otimes \vec{w}$,
is a new object that lives in a higher-dimensional space. It's defined as:

$$
\vec{v} \otimes \vec{w} = (v_1 w_1, v_1 w_2, v_1 w_3, v_2 w_1, v_2 w_2, v_2 w_3)
$$

This is a 6-dimensional vector, where each component is a product of a component from
$\vec{v}$ and a component from $\vec{w}$.

More generally, if $\vec{v}$ is an $n$-dimensional vector and $\vec{w}$ is an $m$-dimensional vector,
their tensor product $\vec{v} \otimes \vec{w}$ is an $nm$-dimensional vector.
\\

Now, let's look at the tensor product of two matrices.
If $A$ is an $n \times m$ matrix and $B$ is a $p \times q$ matrix,
their tensor product $A \otimes B$ is an $np \times mq$ matrix:

$$
A \otimes B = 
\begin{pmatrix} 
a_{11}B & a_{12}B & \cdots & a_{1m}B \\
a_{21}B & a_{22}B & \cdots & a_{2m}B \\
\vdots & \vdots & \ddots & \vdots \\
a_{n1}B & a_{n2}B & \cdots & a_{nm}B
\end{pmatrix}
$$

Each block in this block matrix is a scalar multiple of the matrix $B$.

The above is actually a specialized version of the tensor product called the
\textbf{Kronecker product}, see
\href{https://en.wikipedia.org/wiki/Kronecker_product}{wikipedia: kronecker products}.
\\

In quantum mechanics, the state space of a composite system is
the tensor product of the state spaces of the individual systems.
If $\ket{\psi}$ is a state vector in one system and $\ket{\phi}$ is a state vector in another system,
the state of the composite system is $\ket{\psi} \otimes \ket{\phi}$
the notation $\ket{\psi \phi}$ is often used as a shorthand,
so that the tensor product $\ket{\psi} \otimes \ket{\phi} = \ket{\psi \phi}$,
especially when dealing with composite systems.

The tensor product $\ket{\psi} \otimes \ket{\phi}$ is a formal mathematical operation that combines
two state vectors into a single state vector in a higher-dimensional space.
This is the correct way to describe the state of a composite system in quantum mechanics.

If $\ket{\psi}$ is the state of one particle and $\ket{\phi}$ is the state of
another particle, then $\ket{\psi \phi}$ would typically be understood to mean
$\ket{\psi} \otimes \ket{\phi}$.
\\~\\

The tensor product itself is not a Hermitian form, but it can be used to construct a Hermitian form.

A Hermitian form on a complex vector space $V$ is a map
$h: V \times V \to \mathbb{C}$.
A sesquilinear form then it is a map $h: \overline{V} \times V \to \mathbb{C}$.
\\

The universal property of the tensor product is a key concept that
characterizes the tensor product in terms of a unique mapping property.
The "universal-property definition" of the tensor product of two vector spaces is the following
(see \href{https://en.wikipedia.org/wiki/Tensor_product#Universal_property}{wikipedia: tensor product}):

The tensor product of two vector spaces $V$ and $W$ is a vector space denoted as
$V\otimes W$, together with a bilinear map $\otimes : (v,w) \mapsto v\otimes w$
from $V\times W$ to $V\otimes W$, such that, for every bilinear map
$h : V\times W\to Z$, there is a unique linear map $\tilde{h} :V\otimes W\to Z$,
such that $h= \tilde{h} \circ \otimes$,
that is, $h(v,w) = \tilde{h} (v\otimes w)$ for every $v\in V$ and $w\in W$.


In other words, the tensor product $U \otimes V$ is the "most general" vector space that can be constructed
from $U$ and $V$ using a bilinear map.
Any bilinear map from $U \times V$ to another vector space $W$ factors uniquely through the tensor product
via a linear map from $U \otimes V$ to $W$.
\\

The phrase "up to isomorphism" is a common one in abstract algebra.
It means that two mathematical objects (in this case, vector spaces) are not necessarily identical,
but they are isomorphic, meaning that there exists an isomorphism between them.

An isomorphism between two vector spaces $V$ and $W$ is a bijective (one-to-one and onto)
linear map $f: V \to W$.
If such a map exists, we say that $V$ and $W$ are isomorphic, and we write $V \cong W$.





%%%%%%%%%%%%%%%%%%%%%%%%%%%%%%%%%
\subsection{Practical definition}

Given two vector spaces $V$ and $W$ with basis vectors $e^i$ and $e^j$,
respectively, the tensor product of these spaces, denoted as $V \otimes W$,
is \textbf{a new vector space spanned} by the basis vectors $e^i \otimes e^j$ 
for all possible combinations of $i$ and $j$.

For example, if $V$ is a 2D space with basis vectors $\{e^a, e^b\}$ and
$W$ is a 3D space with basis vectors $\{e^1, e^2, e^3 \}$,
then the tensor product space $V \otimes W$ will be a 6D space with basis vectors:
$$
\{e^a \otimes e^1, e^a \otimes e^2, e^a \otimes e^3, e^b \otimes e^1, e^b \otimes e^2, e^b \otimes e^3\}
$$

To understand this further, let's look at $e^a \otimes e^1$.
Consider the original vector spaces $V$ and $W$.

In the 2D space $V$, $e^a$ might represent the unit vector along the $x$-axis, $(1, 0)$.
In the 3D space $W$, $e^1$ might represent the unit vector along the $x$-axis, $(1, 0, 0)$.
Remember that when we say vectors, we are actually thinking about column vectors.
And it turns out the way to obtain an explicit answer is through the kronecker product again.
In this case, our tensor product can be computed as an \textbf{outer product}.
Remember that the outer product of two vectors $x$ and $y$ is defined as $xy^\dagger$, or
simply $xy^T$, if we are dealing with solely real numbers.
This is in contrast with an inner product which would be $x^\dagger y$ or $x^T y$.

So now we have a great comparison between non-degenerate Hermitian forms and outer product.
(Remember that for a Hermitian form to be an iner product we need the Hermitian form to be
positive definite, that is $(x | y) = \braket{x | y }$ if
 $(x | y)> 0$, see\href{https://en.wikipedia.org/wiki/Inner_product_space}{wikipedia: inner product space}.)

$$
u \otimes v = uv^\dagger
$$
and
$$
\braket{u, v} = \braket{u | v} = u^\dagger v
$$

So when we go on talking about components as in
$$
u \otimes v 
= (u^i e_i) \otimes (v^j b_j)
= u^i v^j e_i \otimes b_j
$$

It means that
$$
(u \otimes v)_{ij} = u^i v^j
$$
That is, the tensor product can be pieced together through normal products between components.

For example,

$$
u \otimes V
= u v^\dagger
=
\begin{pmatrix}
    u_1 \\
    u_2 \\
    \vdots \\
    u_n
\end{pmatrix}
\begin{pmatrix}
    v_1 & v_2 & \dots & v_m
\end{pmatrix}
=
\begin{pmatrix}
u_1 v_1 & u_1 v_2 & \dots  & u_1 v_m \\
u_2 v_1 & u_2 v_2 & \dots  & u_2 v_m \\
\vdots  & \vdots  & \ddots & \vdots \\
u_n v_1 & u_n v_2 & \dots & u_n v_m \\
\end{pmatrix}
$$

In practice, $e^a \otimes e^1$ can be thought of as a 6D unit vector that points in the direction
corresponding to the combination of the x-axis direction from space $V$ and the x-axis direction from space $W$.
It is a basis vector that allows you to describe vectors in the 6D space that have components along this combined direction.

Using what we noted above,
$$
e^a \otimes e^1
= 
\begin{pmatrix}
    1 \\
    0
\end{pmatrix}
\begin{pmatrix}
    1 & 0 & 0
\end{pmatrix}
=
\begin{pmatrix}
u_1 v_1 & u_1 v_2 & u_1 v_3 \\
u_2 v_1 & u_2 v_2 & u_2 v_3 \\
\end{pmatrix}
=
\begin{pmatrix}
(1) (1) & (1) (0) & (1) (0) \\
(0) (1) & (0) (0) & (0) (0) \\
\end{pmatrix}
=
\begin{pmatrix}
1 & 0 & 0 \\
0 & 0 & 0 \\
\end{pmatrix}
$$

Here is where we make our comback to quantum mechanics.
If $\ket{\psi} \in V$ and $\ket{\phi} \ in W$, then a member of the tensor product $V \otimes W$ can be written
as
$$
\ket{\psi}\ket{\phi} = \ket{\psi, \phi} = \ket{\psi\phi} = \ket{\psi} \otimes \ket{\phi}
$$
See
\href{https://en.wikipedia.org/wiki/Bra%E2%80%93ket_notation}{wikipedia: braket notation - composite bras and kets}.
\\

There is one more interesting fact to know!
The tensor product $e^i \otimes e^i$, when $\{e^i\}$ represents an orthonormal basis set,
will always result in a diagonal matrix.
Since it is the product of a matrix with its transpose.


%%%%%%%%%%%%%%%%%%%%%%%%%%%%%%%%%%%%%%%%%%%%%%%%%%%%%%%%%%%%%%%%%%%%%%%%%%%%%%%%%%%%%
\section{Tensor Contraction}


Tensor contractions are used to reduce the rank of a tensor by summing over pairs of indices.
A more useful way to think of tensor contractions is an operations that
can be used to raise or lower indices when the tensor being contracted with
is a metric tensor.
In general relativity, the metric tensor $g_{ij}$ is used to raise or lower indices:
$$
v^i = g^{ij} v_j
$$
$$
v_i = g_{ij} v^j
$$

Here, $g^{ij}$ is the inverse of the metric tensor $g_{ij}$.

We already saw this example when we began talking about tensors and about the inner product,
\ref{chapter:tensors}.
\\

An even simpler example, of tensor contraction would be the raising or lowering of indices using
the metric tensor in a Euclidean space with a Cartesian coordinate system.

In this case, the metric tensor is just the Kronecker delta, denoted as $\delta_{ij}$,
where $i$ and $j$ are the indices.

The Kronecker delta is defined as:
$$
\begin{cases}
    \delta^i_j = 1 & \text{ if } i = j \\
    \delta^i_j  = 0 & \text{ if } i \neq j
\end{cases}
$$

Now, let's say we have a vector $V^i$ (contravariant vector) and we want to find its
covariant counterpart $V_i$.
We can do this by contracting the vector with the metric tensor:

$$
V_i = \delta^i_j  V^j
$$

Here, we sum over the repeated index $j$, which is the essence of tensor contraction.
This operation lowers the index of the vector, transforming it from a contravariant vector
to a covariant vector.
Similarly, if we have a covariant vector $V_i$ and want to find its contravariant counterpart $V^i$,
we can use the inverse metric tensor $\delta^{ij}$
(which in this case is the same as the metric tensor itself) to raise the index:
$$
V^i = \delta_j^i V_j
$$
Again, we sum over the repeated index $j$ to perform the contraction.

And note how $\delta_j^i$ is the inverse of $\delta^j_i$.
\\


The summation over indices is what's shown in the book with the more formal syntax.

Given $T \in T_{s}^{r} (V)$ which has components $T_{i_1, \dots , 1_r }^{j_1 , \dots , js}$
it forms a tensor product of the form
$$ 
V^{*}_{r} \times \dots \times V^{*}_{r} \times V_s \times \dots \times V_s
$$

so that,
$$
T = T_{i_1, \dots , 1_r }^{j_1 , \dots , j_s} e^{i_1} \otimes \dots \otimes e^{i_r} \otimes e_{j_1} \otimes \dots \otimes e_{j_s}
$$

The reason that the tensor components have $r$ products of $V^*$ and $s$ product of $V$, is because
taking the tensor product with their basis vectors will produce a scalar value, and actual component.
And because it allows us to visualize the procedure of a contraction as feeding $e_i$
into the $r$-th slot and $e^i$ into the $(r+s)$-th slot.

Sicne the $r$-th slot corresponds to the basis vector $e^{i_r}$, and the $(r+s)$-th slot corresponds
to the basis vector $e_{j_s}$.
$$
T_{i_1, \dots, i_{r-1}, i}^{j_1, \dots, j_{s-1}, i} =
T_{i_1, \dots, i_r}^{j_1, \dots, j_s} e^{i_1} \otimes \dots \otimes e^{i_{r-1}} \otimes e^i \otimes e_{j_1} \otimes \dots \otimes e_{j_{s-1}} \otimes e_i
$$

Then perform the summation over the repeated index $i$.

Which as we saw in \ref{chapter:tensors}, the $e^i \otimes e_i$ component results in a $\delta^{i}_{i}$
like term, that then makes sure that only a set of matching indices is kept in sync, while
the other indices are used in a summation (one set of indices is kept in sync, the others vary from
0 to $N$ and are then summed at the end).
\section{Integrals in QFT}

\subsection{Kinetic operator and a Green Function} \label{integral:kinetic-operator}

This problem is part of
\href{https://arxiv.org/abs/physics/0609191}{MontePython: Implementing Quantum Monte Carlo using Python}.

There is a kinetic operator defined as
$$
G_K \left(\mathbf{R}^\prime, \mathbf{R}, t\right) =
\frac{1}{\left(2\pi\right)^{3N}} \int e^{-i\mathbf{k}\mathbf{R}^\prime} e^{-Dtk^2} e^{-i\mathbf{k}\mathbf{R}} d\mathbf{k}
$$

This integral then is integrated and the following Green function is obtained
$$
G_K \left(\mathbf{R}^\prime, \mathbf{R}, t\right) =
\frac{1}{\left(4\pi Dt\right)^{3N/2}} e^{-\left(\mathbf{R}-\mathbf{R}^\prime\right) / 4Dt}
$$

And as we saw in \hyperref[jackson:problem-1.2]{Jackson problem 1.2}, the Green function we just got is equivalent to
$\delta\left(\mathbf{R} - \mathbf{R}^\prime\right)$.

But coming back to the integral, the trick is to complete the square!
\begin{align*}
G_K \left(\vect{R}^\prime, \vect{R}, t\right) &=
\frac{1}{\left(2\pi\right)^{3N}} \int e^{-i\vect{k}\vect{R}^\prime} e^{-Dtk^2} e^{-i\vect{k}\vect{R}} d\vect{k} \\
&= \frac{1}{\left(2\pi\right)^{3N}} \int e^{-Dtk^2 + -i\vect{k}\cdot\left(\vect{R}-\vect{R}^\prime\right)} d\mathbf{k}
\end{align*}

This last integral has the form $\int e^{ax^2 + bx} dx$, which can be solved as follows,
\begin{align*}
\int e^{ax^2 + bx} dx &= \int \exp\left\{ a\left(x^2 + \frac{b}{a}x\right) \right\} dx \\
&= \int \exp\left\{ a\left(x^2 + \frac{b}{a}x + \left(\frac{b}{2a}\right)^2 - \left(\frac{b}{2a}\right)^2 \right) \right\} dx \\
&= \int \exp\left\{ a\left(x + \frac{b}{2a}\right)^2 - \frac{b^2}{4a} \right\} dx \\
&= e^{-b^2/4a} \int e^{ a\left(x + \frac{b}{2a}\right)^2 } dx
\end{align*}

Now, let's look at $\int e^{ a\left(x + b/2a\right)^2 } dx$.
First, we need to make a transformation.
Let's define $u = x + b/2a$, then $du = dx$, the measure remains invariant (so do the limits of integration), so
$\int e^{ a\left(x + b/2a\right)^2 } dx = \int e^{ au^2 } du$.
Now, if we define $a = -c$, then 

\begin{align*}
\int e^{ a\left(x + b/2a\right)^2 } dx &= \int e^{ au^2 } du \\
&= \int e^{ -cu^2 } du \\
&= \sqrt{ \frac{\pi}{c} } = \sqrt{ \frac{\pi}{-a} } 
\end{align*}

This whole thing works if $a < 0$ but it turns out that this is also valid if $\re{(a)} \leq 0$ but $a \neq 0$.

Putting everything back together,
$$
\int e^{ax^2 + bx} dx = \sqrt{ \frac{\pi}{-a} } e^{-b^2/4a}
$$

Looking back at our original problem, we can chose $a = -Dt$ and $b = i \vect{r}$, where $\vect{r} = \vect{R}-\vect{R}^\prime$.
and so
\begin{align*}
G_K \left(\vect{R}^\prime, \vect{R}, t\right) &=
\frac{1}{\left(2\pi\right)^{3N}} \int e^{-Dtk^2 + -i\vect{k}\cdot\vect{r}} d\mathbf{k} \\
&= \frac{1}{\left(2\pi\right)^{3N}} \int e^{-Dt \left( k^2 + \frac{i\vect{k}}{Dt}\cdot\vect{r} \right)} d\mathbf{k} \\
&= \frac{1}{\left(2\pi\right)^{3N}} \int e^{-Dt \left( k^2 + \frac{i\vect{r}}{Dt}\cdot\vect{k} + \left(\frac{i\vect{r}}{2Dt}\right)^2 - \left(\frac{i\vect{r}}{2Dt}\right)^2 \right)} d\mathbf{k} \\
&= \frac{1}{\left(2\pi\right)^{3N}} \int e^{-Dt \left( k + \frac{i\vect{r}}{2Dt}\right)^2 - \frac{ \vect{r}^2 }{ 4Dt } } d\mathbf{k} \\
&= \frac{1}{\left(2\pi\right)^{3N}} e^{-\vect{r}^2 /4Dt } \int e^{-Dt \left( k + \frac{i\vect{r}}{2Dt}\right)^2 } d\mathbf{k} \\
\end{align*}

Again, since $\int e^{-ax^2} dx = \sqrt{ \frac{\pi}{a} }$,
then
\begin{align*}
\int e^{-Dt \left( k + \frac{i\vect{r}}{2Dt}\right)^2 } dk &= \sqrt{ \frac{\pi}{Dt} }
\end{align*}

Thus,
\begin{align*}
G_K \left(\vect{R}^\prime, \vect{R}, t\right) &=
\frac{1}{\left(2\pi\right)^{3N}} e^{-\vect{r}^2 /4Dt } \int e^{-Dt \left( k + \frac{i\vect{r}}{2Dt}\right)^2 } d\mathbf{k} \\
&= \frac{1}{\left(2\pi\right)^{3N}} e^{-\vect{r}^2 /4Dt } \left( \sqrt{ \frac{\pi}{Dt} } \right)^{3N} \\
&= \frac{1}{\left(4\pi^2\right)^{3N/2}} \frac{\pi^{3N/2}}{(Dt)^{3N/2}} e^{-\vect{r}^2 /4Dt } \\
&= \frac{1}{\left(4\pi Dt\right)^{3N/2}} e^{-\left(\vect{R} - \vect{R}^\prime\right)^2 /4Dt }
\end{align*}



%%%%%%%%%%%%%%%%%%%%%%%%%%%%%%%%%%%%%%%%%%%%%%%%%%%%%%%%%%%%%%%%%%%%%%%%%%%%%%%%%%%%%%%%%%%%%%%%%%%%%%%%%

\subsection{Klein-gordon Non-relativistic Amplitude} \label{ps-kg:nonrelativistic-propagator}

The following integral comes from Peskin and schroeder Section 2.1
\begin{align}
& \int \frac{d^3p}{(2\pi)^3} \, e^{-i(\vect{p}^2/2m)t} e^{i\vect{p}\cdot(\vect{x}-\vect{x}_0)} \label{kg:nonrelativistic-03} \\
&= \left( \frac{m}{2\pi i t} \right)^{3/2} e^{im (\vect{x}-\vect{x}_0)^2 /2t} \label{kg:nonrelativistic-04}
\end{align}

This integral follows the same procedure as \ref{integral:kinetic-operator}.

Here we have $a = -\frac{it}{2m}$ and $b = i(\vect{x}-\vect{x_0})$.
Since $\re{a} = -t/2m \leq 0$ and $a\neq 0$ when $t\neq 0$ oru solution is valid for $t\neq 0$.
Then

$$
\sqrt{ \frac{\pi}{-a} } = \sqrt{ \frac{2\pi m}{it} }
$$

and
\begin{align*}
-\frac{b^2}{4a} &= \frac{2m \left(\vect{x}-\vect{x}_0\right)^2}{-4it} \\
&= \frac{im \left(\vect{x}-\vect{x}_0\right)^2}{2t}
\end{align*}

So the solution is
\begin{align*}
& \frac{1}{(2\pi)^3} \left( \frac{2\pi m}{it} \right)^{3/2} e^{\frac{im \left(\vect{x}-\vect{x}_0\right)^2}{2t}} \\
&= \left( \frac{m}{2\pi it} \right)^{3/2} e^{im \left(\vect{x}-\vect{x}_0\right)^2 /2t}
\end{align*}







%%%%%%%%%%%%%%%%%%%%%%%%%%%%%%%%%%%%%%%%%%%%%%%%%%%%%%%%%%%%%%%%%%%%%%%%%%%%%%%%%%%%%%%%%%%%%%%%%%%%%%%%%
\section{Tensors}

Working with a spacetime metric signature $(+ - - -)$,
$$
\begin{pmatrix}
    1 & 0  & 0  & 0  \\
    0 & -1 & 0  & 0  \\
    0 & 0  & -1 & 0  \\
    0 & 0  & 0  & -1 \\
\end{pmatrix}
$$

$x^\mu = (x^0, \vect{x})$ and $x_\mu = g_{\mu\nu} x^\nu = (x^0, -\vect{x})$.
To see this explciitly, let's carry out the summations.

When $\mu = 0$, $x_0 = g_{00}x^0 + g_{01}x^1 + g_{02}x^2 + g_{03}x^3 = x^0 + 0 + 0 +0 = x^0$.
For $\mu = 1$, $x_1 = g_{10}x^0 + g_{11}x^1 + g_{12}x^2 + g_{13}x^3 = 0 + x^1 + 0 +0 = -x^0$.
Similarly, when $\mu = 2$, $x_2 = -x^2$, and when $\mu =3$, $x_3 = -x^3$.

The contravariant components of the vector $x^\mu$ are just its components.
The covariant compoents $x_\mu$, defined to be equal to $x_\mu = x^\nu g_{\nu\mu}$, are actually the components
of the associated dual vector.
Why? well, in order to obtain $x_\mu$ we had to use a "non-degenerate hermitian form"
(something that looks like an inner product; though it isn't since this metric is not positive definite).
This makes it a dual vector by definition.
\\

Also note that $g_{\mu\nu}$ forms a \textbf{bilinear form} since $g(x_1, x_2) = g(x_2, x_1)$
(linearity on both sides of the non-degenerate hermitian form).


%%%%%%%%%%%%%%%%%%%%%%%%%%%%%%%%%%%%%%%
\subsection{Inner Product}

Why does $p \cdot x = g_{\mu\nu} p^\mu x^\nu = g^{\mu\nu} p_\mu x_\nu = p_\mu x^\nu = p^\mu x_\nu = p^0 x^0 - \vect{p}\cdot\vect{x}$?

To start, we ought to know that there is the implicit defintion of
$x_\mu = x^\nu g_{\nu\mu}$.
So the covariant vector is defined through the application of a non-degenerate hermitian form
and it constitutes a dual vector.
From there the rest is linear algebra.

The metric tensor defines the geometry of spacetime, including the way distances and angles are measured.
The implied summation, einstein summation, effectively 'weights' the components according to the geometry of spacetime.

$g_{\mu\nu}$ lowers indices, while $g^{\mu\nu}$ raises them.
$\delta^{\mu}_{\nu}$ is a diagonal matrix with ones on the diagonal and zeros elsewhere.
It selects the $\mu$-th component when used in a summation, acting like an identity element.

$g^{\mu\alpha} g_{\alpha\nu} = \delta^{\mu}_{\nu}$

$$
g^{\mu\alpha} g_{\alpha\nu} = \sum_{\alpha = 0} g^{\mu\alpha} g_{\alpha\nu} =
g^{\mu 0} g_{0 \nu} + g^{\mu 1} g_{1\nu} + g^{\mu 2} g_{2\nu} + g^{\mu 3} g_{3\nu}
$$

This is essentially a matrix multiplication.
$$
\begin{pmatrix}
    1 & 0  & 0  & 0  \\
    0 & -1 & 0  & 0  \\
    0 & 0  & -1 & 0  \\
    0 & 0  & 0  & -1 \\
\end{pmatrix}
\cdot
\begin{pmatrix}
    1 & 0  & 0  & 0  \\
    0 & -1 & 0  & 0  \\
    0 & 0  & -1 & 0  \\
    0 & 0  & 0  & -1 \\
\end{pmatrix}
=
\begin{pmatrix}
    1 & 0 & 0 & 0 \\
    0 & 1 & 0 & 0 \\
    0 & 0 & 1 & 0 \\
    0 & 0 & 0 & 1 \\
\end{pmatrix}
$$

Going component by component, take the second row and second column, $\mu=1$ and $\nu=1$.
The element for that entry is given by
$$
g^{1\alpha}g_{\alpha 1} =
g^{1 0} g_{0 1} + g^{1 1} g_{1 1} + g^{1 2} g_{2 1} + g^{1 3} g_{3 1}
= 0 + (-1)(-1) + 0 + 0 = 1
$$


%%%%%%%%%%%%%%%%%%%%%%%%%%%%%%%%%%%%%%%
\subsection{Derivatives}

The examples here come from \cite{tensor-calc}.
The aim here is to see a few examples to get used to the quantities we are going to be seeing later on.
\\

Some useful identities before we carry on:
\begin{enumerate}
    \item
    $$
    \frac{\partial x_p}{\partial x_q} = \delta_{pq}
    $$
\end{enumerate}



\textbf{Example 1}

If $a_{ij}$ are contants, to compute the partial derivative
$$
\frac{\partial}{\partial x_k} a_{ij} x_i x_j
$$

We'll firs write out the sums explciitly.
\begin{align*}
a_{ij} x_i x_j &= \sum_{i,j} a_{ij} x_i x_j
\end{align*}

Someone once happened to noticed that the sum could be broken into parts depending on when $x_{\{i,j\}} = x_k$, where $x_k$
is the one we want to differentiate with respect to,
\begin{align*}
\sum_{i,j} a_{ij} x_i x_j &=
    \sum_{i \neq k, j \neq k} a_{ij} x_i x_j +
    \sum_{i, j \neq k} a_{ij} x_i x_j + 
    \sum_{i \neq k, j} a_{ij} x_i x_j +
    \sum_{i=k, j=k} a_{ij} x_i x_j \\
&= C + \sum_{j \neq k} \left( a_{kj} x_j \right) x_k +
    \sum_{i \neq k} \left( a_{ik} x_i \right) x_k +
    a_{kk} x_k x_k \\
    &= C + \sum_{j \neq k} \left( a_{kj} x_j \right) x_k +
    \sum_{i \neq k} \left( a_{ik} x_i \right) x_k +
    a_{kk} (x_k)^2 \\
\end{align*}

We wrote the first sum, where no $x$s are equal to $x_k$ as $C$, because as you may guess this is going to go away with
any derivative.

Back to our original task,
\begin{align*}
\frac{\partial}{\partial x_k} a_{ij} x_i x_j &=
    0 + \sum_{j \neq k} a_{kj} x_j +
    \sum_{i \neq k} a_{ik} x_i +
    2 a_{kk} x_k \\
&= \sum_{j} a_{kj} x_j + \sum_{i} a_{ik} x_i \\
&= a_{kj} x_j + a_{ik} x_i \\
&= a_{ki} x_i + a_{ik} x_i \\
&= (a_{ik} + a_{ki}) x_i
\end{align*}



\textbf{Example 2}

Another way to work out
$$
\frac{\partial}{\partial x_k} a_{ij} x_i x_j
$$
goes as follows.

\begin{align*}
\frac{\partial}{\partial x_k} a_{ij} x_i x_j &= a_{ij} \frac{\partial}{\partial x_k} x_i x_j \\
&= a_{ij} \left( \frac{\partial x_i}{\partial x_k} x_j + x_i \frac{\partial x_j}{\partial x_k} \right) \\
&= a_{ij} \left( \delta_{ik} x_j + x_i \delta_{jk} \right) \\
&= a_{kj} x_j + a_{ik} x_i \\
&= (a_{ik} + a_{ki}) x_i
\end{align*}


\textbf{Example 3}

If $a_{ij} = a_{ji}$ are constants, compute
$$
\frac{\partial^2}{\partial x_k \partial x_l} a_{ij} x_i x_j
$$

Let's start with our previous result,
\begin{align*}
\frac{\partial^2}{\partial x_k \partial x_l} a_{ij} x_i x_j &=
    \frac{\partial}{\partial x_k} \left[ \frac{\partial}{\partial x_l} a_{ij} x_i x_j \right] \\
&= \frac{\partial}{\partial x_k} (a_{il} + a_{li}) x_i \\
&= 2 a_{li} \frac{\partial x_i}{\partial x_k} \\
&= 2 a_{li} \delta_{ik} \\
&= 2 a_{lk}
\end{align*}



%%%%%%%%%%%%%%%%%%%%%%%%%%%%%%%%%%%%%%%%%%%%%%%%%%%%%%%%%%%%%%%%%%%%%%%%%%
\subsection{Changing Variables: the Jacobian}


Promise we are not doing the conventional thing where we go from scratch but there is one more pattern we have to covered
now in order for some other equations to make sense and this is the Jacobian.


There is a result from calculus that is presented a bit like this (the following presentation comes from
\cite{jacobian}):


If we define the Jacobian of a transformation $x = g(u,v)$, $y = h(u,v)$ as
\begin{align*}
\frac{ \partial (x,y) }{ \partial (u,v) } &= 
\begin{vmatrix}
\dfrac{ \partial x }{ \partial u } & \dfrac{ \partial x }{ \partial v } \\ \addlinespace
\dfrac{ \partial y }{ \partial u } & \dfrac{ \partial y }{ \partial v }
\end{vmatrix} \\
&= \frac{ \partial x }{ \partial u } \frac{ \partial y }{ \partial v } -
    \frac{ \partial x }{ \partial v }\frac{ \partial y }{ \partial u }
\end{align*}

Then if we want to integrate $f(x)$ over some surface $R$, this region will become some other surface $S$
following our transformation and the integral can be computed as follows:
$$
\int \int_{R} f(x,y) dA
= \int \int_{S} f\left( g(u,v), h(u,v) \right) \Biggl| \frac{ \partial (x,y) }{ \partial (u,v) } \Biggl| d\overline{A}
$$

For example, when we go from Euclidean to polar coordinates, the transformation we use is $x = r\cos\theta$
and $y = r\sin\theta$ and the Jacobian is $r$ so
$dA = \left| \frac{\partial (x, y)}{\partial (r, \theta)} \right| dx dy = r dr d\theta$:

\begin{align*}
\frac{\partial (x, y)}{\partial (r, \theta)} &= 
\begin{vmatrix}
    \dfrac{ \partial x }{ \partial r } & \dfrac{ \partial x }{ \partial \theta } \\ \addlinespace
    \dfrac{ \partial y }{ \partial r } & \dfrac{ \partial y }{ \partial \theta }
\end{vmatrix} \\
& = \begin{vmatrix}
    \cos\theta & -r\sin\theta \\ \addlinespace
    \sin\theta & r\cos\theta
\end{vmatrix} \\
&= r \cos^2 \theta + r \sin^2 \theta \\
&= r
\end{align*}



In 3-dimensions this extends to
$$
\int \int \int_{V} f(x,y,z) dV
= \int \int \int_{B} f\left( g(u,v,w), h(u,v,w), k(u,v,w) \right) \Biggl| \frac{ \partial (x,y,z) }{ \partial (u,v,w) } \Biggl| d\overline{V}
$$

For the case of the transformation of Euclidean coordinates to spherical coordinates, the transformation of variables is
$x = r \sin\phi \cos\theta $, $y = r \sin\phi \sin\theta$, and $z = r\cos\phi$.
In that case the Jacobian is $r^2 \sin\phi$ and thus
$dV = \left| \frac{\partial (x, y, z)}{\partial (r, \theta, \phi)} \right| dx dy dz = r^2 \sin\phi dr d\phi d\theta$.


Besides being very widely used, the reason we mentioned it is because this is also that is very well used in
\cite{gifted-qft} to introduce the notation we see so much in QFT.
Specifically, the Jacobian as we have been writing it, can also be seen in expressions such as

$$
\overline{a}^{\mu} = \left( \frac{\partial \overline{x}^\mu}{\partial x^\nu} \right) a^\nu
$$

$$
\frac{\partial \phi}{\partial \overline{x}^\mu} =
    \left( \frac{\partial x^\nu}{\partial \overline{x}^\mu} \right) \frac{\partial \phi}{\partial x^\nu}
$$

The determinant of the Jacobian matrix is used specifically when you are dealing with volume transformations
(like integrating over a volume or changing the measure of integration) because it scales the volume elements according
to how the coordinate transformation stretches or compresses the space.

When transforming vectors, however, you're not scaling a volume but rather reorienting or rescaling each component
of the vector according to how the coordinate axes themselves have changed.
This is why the transformation involves a sum over the product of vector components and the appropriate Jacobian elements:

\begin{itemize}
\item Contravariant vectors (standard vectors): The components of these vectors are transformed by multiplying with the
    Jacobian matrix $\partial \overline{x}^\mu / \partial x^\nu$. Here, each new vector component is a linear combination
    of old components weighted by how much each new coordinate axis changes with respect to each old axis.
    \begin{itemize}
        \item These objects, such as displacement vectors, inherently "follow" the coordinate axes. As the axes stretch
        or rotate, so do the components of the vector. 
    \end{itemize}
\item Covariant vectors (gradients and one-forms): These vectors transform using the inverse of the Jacobian matrix
    $\partial x^\nu / \partial \overline{x}^\mu$. This reflects how changes in the old coordinates map to the new
    coordinates, suitable for objects that naturally pair with vectors, like gradients.
    \begin{itemize}
        \item They measure rates of change along these coordinates.
        \item A larger coordinate stretch means a smaller gradient in that direction to maintain the same rate of change.
    \end{itemize}
\end{itemize}




%%%%%%%%%%%%%%%%%%%%%%%%%%%%%%%%%%%%%%%%%%%%%%%%%%%%%%%
\subsection{From Tensors to Matrices}

Tensors can be expressed as matrices only once a basis has been chosen.
the example being given is with the angular momentum operator $L_z$ on the spherical harmonics $H_l (\mathbb{R}^3)$.

It says that $H_l (\mathbb{R}^3)$ is the set of all linear functions on $\mathbb{R}^3$ that

$$
\{ r Y_{m}^{l} \}_{ -1 \leq m \leq 1 } =
\Big\{ \frac{1}{\sqrt{2}} (x +iy), \, z, \, - \frac{1}{\sqrt{2}} (x -iy) \Big\}
$$

and $\{ x, y, z \}$ are both basis for the space.

Then having the angular momentum operator $L_z = -i ( x\partial_y - y \partial_x )$ on this space we got
$$
\frac{1}{\sqrt{2}} L_z (x + iy) = \frac{1}{\sqrt{2}} (x + iy)
$$
$$
L_z (z) = 0
$$
$$
-\frac{1}{\sqrt{2}} L_z (x -iy) = \frac{1}{\sqrt{2}} (x -iy)
$$

which implies that in the spherical harmonics basis

$$
L_z = 
\begin{pmatrix}
1 & 0 & 0 \\
0 & 0 & 0 \\
0 & 0 & -1
\end{pmatrix}
$$

Could you tell me how that tensor representation of $L_z$ is implies by saying that
$T(e_i) = \sum^{n}_{j=1} T_{i}^{j} e_j$ ?


$$
\frac{1}{\sqrt{2}} L_z (x + iy)
\rightarrow
L_z \left( \frac{1}{\sqrt{2}} (x + iy) \right)
$$

$$
L_z (z) = 0
$$

$$
-\frac{1}{\sqrt{2}} L_z (x -iy)
\rightarrow
L_z \left( -\frac{1}{\sqrt{2}} (x + iy) \right)
$$

Since

$$
\{ r Y_{m}^{l} \}_{ -1 \leq m \leq 1 } =
\Big\{ \frac{1}{\sqrt{2}} (x +iy), \, z, \, - \frac{1}{\sqrt{2}} (x -iy) \Big\}
$$

then we can say that $e_1 = \frac{1}{\sqrt{2}} (x +iy)$, $e_2 = z$, and $e_3 = -\frac{1}{\sqrt{2}} (x -iy)$.

Meaning that $L_z (e_1) = e_1 = 1\cdot e_1 + 0\cdot e_2 + 0\cdot e_3$,
$L_z (e_2) = 0 = 0\cdot e_1 + 0\cdot e_2 + 0\cdot e_3$, and
$l_z (e_3) = -e_3 = 0\cdot e_1 + 0\cdot e_2 - 1\cdot e_3$.
From here, we can then adopt write this as a matrix

$$
T_{i}^{j}
=
L_z
=
\begin{pmatrix}
1 & 0 & 0 \\
0 & 0 & 0 \\
0 & 0 & -1
\end{pmatrix}
$$

The columns of a matrix representation of a linear operator correspond to the action of the operator on the basis vectors.
However, due to the specific diagonal form of the matrix in this example, both the row and column perspectives
give the same result, which can indeed be confusing.

The first column represents $L_z(e_1)$,
the second column represents $L_z(e_2)$,
the third column represents $L_z(e_3)$.

In physics speak this would be read as the wave functions
$\frac{1}{\sqrt{2}} (x +iy)$, $z$, and $- \frac{1}{\sqrt{2}} (x -iy)$
have eigenvalues 1, 0, and -1, respectively.




%%%%%%%%%%%%%%%%%%%%%%%%%%%%%%%%%%%%%%%%%%%%%%%%%%%%%%%%%%%%%%%%%%%%%%%%%%%%%%%%%%%%%%
\section{Now Let's Talk Tensors}

%%%%%%%%%%%%%%%%%%%%%%%%%%%%%%%%%%%%%%%%%%%%%%%%%%%%%%%%%%%%%%%%%%%%%%%%%%%%%%%%%%%%%%
\section{Transformation Matrices}

The previous section gave us a good place to begin talking about tensors.
Now let's actually talk about them.

Tensors have been historically defined based on how they transform.
Think of the conventional axis rotation formula in 2D.
If we rotate the axis by an angle $\theta$, then

$$
x^\prime = x \cos\theta - y \sin\theta
$$
and
$$
y^\prime = x \sin\theta + y \cos\theta
$$

As we saw in the previous section,
where we derived the matrix representation of the angular momentum operator,
the columns of a matrix representation of a multilinear map (tensor) correspond to the action of
the map (tensor) on the basis vectors.
Which is most commonly seen as

$$
\begin{pmatrix}
x^\prime \\
y^\prime
\end{pmatrix}
=
\begin{pmatrix}
\cos\theta & - \sin\theta \\
\sin\theta & \cos\theta
\end{pmatrix}
\begin{pmatrix}
x \\
y
\end{pmatrix}
$$

If we envision $x$ and $y$, and $x^\prime$ and $y^\prime$, and some basis sets,
such that $x \rightarrow e_1$ and $y \rightarrow e_2$ (similarly for the primed components).
We can rewrite the above as
$$
e_{1^\prime} = e_1 \cos\theta - e_2 \sin\theta
$$
and
$$
e_{2^\prime} = e_1 \sin\theta + e_2 \cos\theta
$$

Meaning that our problem becomes,
$$
\begin{pmatrix}
e_{1^\prime} \\
e_{2^\prime}
\end{pmatrix}
=
\begin{pmatrix}
\cos\theta & - \sin\theta \\
\sin\theta & \cos\theta
\end{pmatrix}
\begin{pmatrix}
e_1 \\
e_2
\end{pmatrix}
$$

Following the Einstein summation convention, we could write

$$
e_{i^\prime} = A_{i^\prime}^{j} e_j
$$

which also means that the original basis can be obtained as follows
$$
e_{i} = A_{i}^{j^\prime} e_{j^\prime}
$$

It is interesting to now quickly pursue the exercise of expressing our basis in terms of itself
as it will show us some properties about these numbers $A_{i^\prime}^{j}$ and $A_{i}^{j^\prime}$.


Combining the previous two expressions we get,
$$
e_{i} = A_{i}^{j^\prime} e_{j^\prime}
= A_{i}^{j^\prime} A_{j^\prime}^{k} e_k
$$

From here we can see that
$$
A_{i}^{j^\prime} A_{j^\prime}^{k} = \delta_{i}^{k}
$$

If we consider the same example but with the primed-basis.

$$
e_{i^\prime} = A_{i^\prime}^{j} e_j 
= A_{i^\prime}^{j} A_{j}^{k^\prime} e_{k^\prime}
$$

Meaning that

$$
A_{i^\prime}^{j} A_{j}^{k^\prime} = \delta_{i^\prime}^{k^\prime}
$$

So the different $A$s are inverse of one another.

It is also worth noting that These $A$s are not tensor components since their indices refer to different basis.
In expressions such as $A_{i^{\prime}}^{j} A_{j}^{k^{\prime}}$, the indices $i^{\prime}$ and $k^{\prime}$ refer
to the primed basis, while the index $j$ refers to the unprimed basis.
This mixing of basis sets in a single expression is why the numbers $A_{i^{\prime}}^{j}$ and $A_{j}^{k^{\prime}}$
are not considered tensor components.

The transformation matrices $A_{i^{\prime}}^{j}$ and $A_{j}^{k^{\prime}}$ are not tensor components themselves,
but rather express the relationship between components in different bases.
The fact that these matrices are inverses of each other, which is a consequence of the transformation law for
tensor components.

Another useful thing to note is that the transformation for dual vectors looks the same but with uppoer indices.

$$
e^{i^\prime} = A_{j}^{i^\prime} e^j
$$
and
$$
e^{i} = A_{j^\prime}^{i} e^{j^\prime}
$$

With this info at hand, we can guesstimate the matrix transformation laws for tensors,

$$
A =
\begin{pmatrix}
A_{1}^{1^\prime} & A_{2}^{1^\prime} & \dots  & A_{n}^{1^\prime} \\
A_{1}^{2^\prime} & A_{2}^{2^\prime} & \dots  & A_{n}^{2^\prime} \\
\vdots           & \vdots           & \ddots & \vdots           \\
A_{1}^{n^\prime} & A_{2}^{n^\prime} & \dots  & A_{n}^{n^\prime}
\end{pmatrix}
$$

Its inverse is
$$
A^{-1}
=
\begin{pmatrix}
A_{1^\prime}^{1} & A_{2^\prime}^{1} & \dots  & A_{n^\prime}^{1} \\
A_{1^\prime}^{2} & A_{2^\prime}^{2} & \dots  & A_{n^\prime}^{2} \\
\vdots           & \vdots           & \ddots & \vdots           \\
A_{1^\prime}^{n} & A_{2^\prime}^{n} & \dots  & A_{n^\prime}^{n} \\
\end{pmatrix}
$$

These are unitary matrices
$$
A A^{-1} = A^{-1} A = I
$$

Equating the above with our 2D rotation example,
$$
A^{-1}
=
\begin{pmatrix}
\cos\theta & - \sin\theta \\
\sin\theta & \cos\theta
\end{pmatrix}
$$


%%%%%%%%%%%%%%%%%%%%%%%%%%%%%%%%%%%%%%%%%%%%%%%%%%%%%%%%%%%%%%%%%%%%%%%%%%%%%%%%%%%%%%
\section{Contravariant and Covariant Vectors}

One fun fact about vectors is that they can be seen as tensors of type (0,1), that is, they feed
on 1 dual vector.
A vector can be viewed as a tensor of type (0, 1) because it transforms according to the contravariant
transformation law under a change of coordinates. The key idea is that a vector represents a directed
quantity that "eats" or operates on dual vectors (covectors or linear functionals) to produce scalars.

This is mathematically said as,
$$
v^{i^\prime} = A^{i^\prime}_{j} v^j
$$
As we saw above, this is the transformation for dual vectors - expressions with contravariant indices.

And as we saw above, these transform according to the Jacobian.

A dual vector in turn transforms as
$$
f_{i^\prime} = A_{i^\prime}^{j} v_j
$$

Note how the transformation are the contrary to what we saw for the basis vectors!
Where as the components of the dual vector transform in the same way as the basis.

In matrix notation we have contravariant vectors transform as
$$
v_{\mathcal{B}^\prime} = A v_{\mathcal{B}}
$$

and covariant ones as
$$
f_{\mathcal{B}^\prime} = \left( A^{-1} \right)^T f_{\mathcal{B}}
$$

But remember that basis and dual basis transform in the opposite ways!
\\


Now, in quantum mechanics, if we consider an orthonormal basis ${\ket{e_{i^{\prime}}}}$ in the same Hilbert space.
And assume that both ${\ket{e_i}}$ and ${\ket{e_{i^{\prime}}}}$ are complete bases,
so there must be a unitary transformation $U$ that relates them.
We would write, $\ket{e_{i^{\prime}}} = U \ket{e_i}$ for this.
Consecuently, the dual basis transforms as follows, $\bra{ e_{i^\prime} } = \bra{e_i} U^\dagger $.

Here, $U$ plays the role of $A^{-1}$, or $A_{i^{\prime}}^{j}$ in the tensor formalism.

Following the previous approach, we get the following results for transformations into a primed basis set,
$$
\ket{e_{i^{\prime}}} = U \ket{e_i}
$$
and
$$
\bra{ e_{i^\prime} } = \bra{e_i} U^\dagger
$$

The matrix elements of $U$ in the $\ket{e_i}$ basis are given by:
$U_{i^{\prime}}^{j} = \bra{e_{i^{\prime}}} U \ket{e_j}$.
These matrix elements play the same role as the $A_{i^{\prime}}^{j}$ in your classical expression.

To see the analogy, let's use the completeness relation in the $\ket{e_i}$ basis, we can write
$$
\ket{e_{i^{\prime}}} = U \ket{e_i}
= \sum_j U \ket{e_j} \braket{e_j | e_i}
= \sum_j U \ket{e_j} \delta_{ij}
= \sum_j U_{i^{\prime}}^{j} \ket{e_j}
$$
This is the quantum analog of your expression $e_{i^{\prime}} = A_{i^{\prime}}^{j} e_j$.

Furthermore, we can use the completeness relation in the $\ket{e_{i^{\prime}}}$ basis to write:
$$
\ket{e_j} = \sum_{k^{\prime}} \ket{e_{k^{\prime}}} \braket{ e_{k^{\prime}} | e_j} = \sum_{k^{\prime}} (U^{\dagger})_{j}^{k^{\prime}} \ket{e_{k^{\prime}}}
$$
Substituting this into the previous expression, we get:
$$
\ket{e_{i^{\prime}}} = \sum_j U_{i^{\prime}}^{j} \sum_{k^{\prime}} (U^{\dagger})_{j}^{k^{\prime}} \ket{e{k^{\prime}}} = \sum_{k^{\prime}} (\sum_j U_{i^{\prime}}^{j} (U^{\dagger})_{j}^{k^{\prime}}) \ket{e_{k^{\prime}}}
$$
This is the quantum analog of your expression $e_{i^{\prime}} = A_{i^{\prime}}^{j} A_{j}^{k^{\prime}} e_{k^{\prime}}$.
\\


%%%%%%% notes from chat

If a state vector $\ket{\psi}$ is expressed in the $\{\ket{e_i}\}$ basis as:

$$
\ket{\psi} = \sum_i c_i \ket{e_i}
$$

then in the $\{\ket{e_{i^{\prime}}}\}$ basis, it is expressed as:

$$
\ket{\psi} = \sum_{i^{\prime}} c_{i^{\prime}} \ket{e_{i^{\prime}}}
$$

where $c^\prime = A c$.
Note that we are not doing the component notation because the above expression is seldom expressed in terms of contravariant
or covariant components, so it is easy to make sense of it this way.

So, the components of the state vector transform with the matrix $A$, while the basis vectors transform with the
matrix $A^{-1}$, just as in the tensor formalism.



%%%%%%%%%%%%%%%%%%%%%%%%%%%%%%%%%%%%%%%%%%%%%%%%%%%%%%%%%%%%%%%%%%%%%%%%%%%%%%%%%%%%%%
\subsection{Pauli Matrices}

The Pauli matrices actually form a basis for $2 \times 2$ Hermitian matrices.
Let's see how: a generic $2 \times 2$ Hermitian matrix would look like

$$
\begin{pmatrix}
a & b\\
c & d
\end{pmatrix}
=
\begin{pmatrix}
a^* & c^* \\
b^* & d^*
\end{pmatrix}
$$

Meaning that $a, b \in \mathbb{R}$ and that $b = c^*$.
So we could write our Hermitian matrix as

$$
\begin{pmatrix}
a      & x - iy \\
x + iy & d
\end{pmatrix}
$$

Another requirement that is imposed on $H_2 (\mathbb{C})$ is that they must be unitary, that is $U^\dagger U = I$, so

\begin{align*}
\begin{pmatrix}
a      & x - iy \\
x + iy & d
\end{pmatrix}
\begin{pmatrix}
a      & x - iy \\
x + iy & d
\end{pmatrix}
&=
\begin{pmatrix}
a^2 + x^2 - y^2         & a(x - iy) + b(x - iy) \\
a(x + iy) + b(x + iy) & d^2 + x^2 - y^2
\end{pmatrix}
\\
&=
\begin{pmatrix}
1 & 0 \\
0 & 1
\end{pmatrix}
\end{align*}

One way to make this happen is to set $a = -d$, this will also have the happy side effect of meeting another requirement
for the Pauli matrices to be traceless.
So now we have our candidate bases that ought to look like

$$
\begin{pmatrix}
a      & x - iy \\
x + iy & -a
\end{pmatrix}
$$

This then means that our requirement for our matrix $H_2 (\mathbb{C})$ to be unitary simplifies to
$$
\begin{pmatrix}
a^2 + x^2 - y^2 & 0                \\
0               & a^2 + x^2 - y^2
\end{pmatrix}
=
\begin{pmatrix}
1 & 0 \\
0 & 1
\end{pmatrix}
$$

Here is where our 3 Pauli matrices comes!

$$
\sigma_1 = \sigma_x =  
\begin{pmatrix}
1 & 0 \\
0 & 1
\end{pmatrix}
$$
comes from assuming that $x=1$, so $a^2 + x^2 - y^2 = 1$ can be met by assuming that $a = y = 0$.

$$
\sigma_2 = \sigma_y =  
\begin{pmatrix}
0 & -i \\
i & 0
\end{pmatrix}
$$

comes from assuming that $y=1$, so $a^2 + x^2 - y^2 = 1$ can be met by assuming that $a = x = 0$.

Finally,
$$
\sigma_3 = \sigma_z =  
\begin{pmatrix}
1 & 0 \\
0 & -1
\end{pmatrix}
$$

Where we took $a=1$ and $x = y = 0$.





%%%%%%%%%%%%%%%%%%%%%%%%%%%%%%%%%%%%%%%%%%%%%%%%%%%%%%%%%%%%%%%%%%%%%%%%%%%%%%%%%
\section{Frobenius Method}

\subsection{A Generic Example}

The following example comes from
\href{https://en.wikipedia.org/wiki/Frobenius_method}{wikipedia: frobenius method}:

$$
z^2 f^{\prime\prime} - zf^\prime + (1-z)f = 0
$$

We assume the solution has the form,
$$
f = \sum_{k=0} A_k z^{k+r}
$$

Plugging the above series into the ODE and organizing like terms we land in
$$
\sum_k \left[ (k+r)(k+r-1) - (k+r) + (1-z) \right] A_k z^{k+r} = 0
$$
which can be further simplified to

$$
\sum_k \left[ (k+r)(k+r-1) - (k+r) + 1 \right] A_k z^{k+r}
    - \sum_k A_k z^{k+r+1} 
= 0
$$

The indicial polynomial can then be obtained by considering the lowest power of $z$, which is $z^r$,
$$
(r(r-1) + 1) A_0 z^r
$$

We now set $z=0$ and obtain
$$
r(r-1) + 1 = r^2 -r + 1 = (r-1)^2 = 0
$$

Meaning that $r=1$ (double root).

We now need to find the recurrence relation by plugging $f = \sum_{k=0} A_k z^{k+1}$ back into our original proble,
equate coefficients of like powers of $z$, and find a recurrence relation for $A_k$.

We now have,
$$
f^\prime = \sum_{k=0} A_k (k+1) z^k
$$
and
$$
f^{\prime\prime} = \sum_{k=0} A_k k(k+1) z^{k-1}
$$

And our ODE now becomes,
\begin{align*}
& z^2 \sum_{k} A_k k(k+1) z^{k-1}
    - z \sum_{k} A_k (k+1) z^k
    + (1-z) \sum_{k} A_k z^{k+1} \\
&= \sum_{k} A_k k(k+1) z^{k+1}
    - \sum_{k} A_k (k+1) z^{k+1}
    + \sum_{k} A_k z^{k+1}
    - \sum_{k} A_k z^{k+2} \\
&= \sum_k \left[ k(k+1) - (k+1) + 1 \right] A_k z^{k+1} - \sum_{k} A_k z^{k+2} \\
&= \sum_k k^2 A_k z^{k+1} - \sum_{k} A_k z^{k+2}
\end{align*}

To find a recurrence relation we need to align the sums.
So we can write
\begin{align*}
& \sum_{k=0} k^2 A_k z^{k+1} - \sum_{k=0} A_k z^{k+2} \\
&= \sum_{k=0} k^2 A_k z^{k+1} - \sum_{k=1} A_{k-1} z^{k+1}
\end{align*}

We could do a similar shift for the first sum but since the first term $k=0$ has a $k^2$ term then we can ignore it.
Thus,
$$
\sum_{k=1} \left( k^2 A_k - A_{k-1} \right) z^{k+1} = 0
$$

Since this sum must be zero for all $z$, the coefficients must satisfy,
$$
k^2 A_k - A_{k-1} = 0
$$
or
$$
A_k = \frac{ A_{k-1} }{k^2} 
$$

This recurrence realtionship now gives us all coefficients in terms of $A_0$.
The first couple coefficients are,
$$
A_1 = \frac{ A_0 }{1^2} = A_0
$$

$$
A_2 = \frac{ A_1 }{2^2} = \frac{ A_0 }{2^2}
$$

$$
A_3 = \frac{ A_2 }{3^2} = \frac{ A_0 }{2^2 3^2}
$$

This means that our general solution looks like
$$
f(z) = \sum_{k=0} A_k z^{k+1}
= A_0 \sum_{k=0} \frac{ z^{k+1} }{k!^2}
$$



%%%%%%%%%%%%%%%%%%%%%%%%%%%%%%%%%%%%%%%%%%%%%%%%%%%%%%%%%%%%%%%%%%%

\subsection{Legendre Equation}

Using the nomenclature of Arfken,
$$
\mathcal{L} y(x) =
-(1-x^2) y^{\prime\prime} + 2x y^\prime = \lambda y
$$

Which can also be expressed as
$$
(1-x^2) y^{\prime\prime} - 2x y^\prime + \lambda y = 0
$$

If we assume the solution is of the form
$$
y(x) =
\sum_{j=0} a_j x^{j+s}
$$

then the Legendre equation can be written as

\begin{align*}
& (1-x^2) \sum_{j} a_j (j+s)(j+s-1) x^{j+s-2}
    - 2x \sum_{j} a_j (j+s) x^{j+s-1}
    + \lambda \sum_{j} a_j x^{j+s} \\
&= \sum_{j} a_j (j+s)(j+s-1) x^{j+s-2}
    - \sum_{j} a_j (j+s)(j+s-1) x^{j+s}
    -2 \sum_{j} a_j (j+s) x^{j+s}
    + \lambda \sum_{j} x^{j+s} \\
&= \sum_{j} a_j (j+s)(j+s-1) x^{j+s-2}
    + \sum_j \left[ -(j+s)(j+s-1) -2(j+s) + \lambda \right] a_j x^{s+j}
\end{align*}

The indicial polynomial then comes from considering the lowest power of $x$, which is the first sum
$$
a_j (j+s)(j+s-1) x^{j+s-2}
$$
which for an arbitrary $x$ when $j=0$, it becomes
$$
s(s-1) = 0
$$
whose solutions are $s=0$ and $s=1$.

In this case, the solution will be a linear combination,

$$
y(x) = y_1 (x) + y_2 (x)
= \sum_j a_j x^{j} + \sum_j b_j x^{j+1}
$$

Let's start with the $s=0$ case then.
Plugging $y_1 (x)$ back into our ODE we have

\begin{align*}
& (1-x^2) \sum_j a_j j(j-1) x^{j-2}
    - 2x \sum_j a_j j x^{j-1}
    + \lambda \sum_j a_j x^j \\
&= \sum_j a_j j(j-1) x^{j-2}
    - \sum_j a_j j(j-1) x^{j}
    - 2 \sum_j a_j j x^{j}
    + \lambda \sum_j a_j x^j \\
&= \sum_j a_j j(j-1) x^{j-2}
    + \sum_j \left[ -j(j-1) -2j + \lambda \right] a_j x^j
\end{align*}

Now, to align the powers of $x^{j}$.
First, note that the first two terms in
$\sum_j a_j j(j-1) x^{j-2}$ are $0$.
So we might as well start at $j=2$
$$
\sum_{j=2} a_j j(j-1) x^{j-2}
$$

If we shift it back to $0$, we can align the powers of $x$.
$$
\sum_{j=2} a_j j(j-1) x^{j-2} =
\sum_{j=0} a_{j+2} (j+2)(j+1) x^{j}
$$

This gives us,
$$
\sum_{j=0} a_{j+2} (j+2)(j+1) x^{j}
+ \sum_j \left[ -j(j-1) -2j + \lambda \right] a_j x^j
= 0
$$

Now, we can equate coefficients,
$$
a_{j+2} (j+2)(j+1)
    + \left[ -j(j-1) -2j + \lambda \right] a_j = 0
$$
or
$$
a_{j+2} = \frac{ j(j-1) + 2j - \lambda }{ (j+1)(j+2) } a_j
= \frac{ j(j+1) - \lambda  }{ (j+1)(j+2) } a_j
$$
\\~\\

In the $s=1$ case, if we plug in $y_2$ into our ODE,

\begin{align*}
& (1-x^2) \sum_j a_j j(j+1) x^{j-1}
    - 2x \sum_j a_j (j+1) x^{j}
    + \lambda \sum_j a_j x^{j+1} \\
&= \sum_j a_j j(j+1) x^{j-1}
    - \sum_j a_j j(j+1) x^{j+1}
    - 2 \sum_j a_j (j+1) x^{j+1}
    + \lambda \sum_j a_j x^{j+1} \\
&= \sum_j a_j j(j+1) x^{j-1}
    + \sum_j \left[ -j(j+1) -2(j+1) + \lambda \right] a_j x^{j+1}
\end{align*}

Following the procedure we have been building: look for a series with redundant terms, and then try to shift it to match
the other powers of $x$.
Here we have
$$
\sum_{j=0} a_j j(j+1) x^{j-1}
$$
that can be written as
$$
\sum_{j=1} a_j j(j+1) x^{j-1}
$$
Since the term when $j=0$ equals to 0.
And we can conveniently shift this series back one to match the $x^j$ terms of the other sum:

$$
\sum_{j=1} a_j j(j+1) x^{j-1}
=
\sum_{j=0} a_{j+1} (j+1)(j+2) x^{j}
$$

At this point we have,

$$
\sum_{j=0} a_{j+1} (j+1)(j+2) x^{j}
    + \sum_j \left[ -j(j+1) -2(j+1) + \lambda \right] a_j x^{j+1}
$$

We need to do another shift and doing it on the first sum seems like the simplest approach.
If we do so, then we can arrive at

$$
\sum_{j=-1} a_{j+2} (j+2)(j+3) x^{j+1}
    + \sum_j \left[ -j(j+1) -2(j+1) + \lambda \right] a_j x^{j+1}
$$

But let's separate the first term of the first sum.
When $j=-1$, we get $a_1 (1)(2) x^0 = 2a_1$.
Thus,

$$
2 a_1 +
    \sum_{j=0} a_{j+2} (j+2)(j+3) x^{j+1}
    + \sum_j \left[ -j(j+1) -2(j+1) + \lambda \right] a_j x^{j+1} = 0
$$

And this equation must hold for all $x$, so the coefficients of each power of $x$ must be zero.
This means that $a_1 must be 0$.
And we are left with
$$
a_{j+2} (j+2)(j+3)
    + \left( -j(j+1) -2(j+1) + \lambda \right) a_j = 0
$$

Which simplifies to
$$
a_{j+2} = \frac{ j^2 + 3j + 2 - \lambda }{ (j+2)(j+3) } a_j
= \frac{ (j+1)(j+2) - \lambda }{ (j+2)(j+3) } a_j
$$

In order to obtain the eignevalues we need to go back to the beginning.
If you remember, this is how the problem was initially framed:
$$
\mathcal{L} y(x) =
-(1-x^2) y^{\prime\prime} + 2x y^\prime = \lambda y
$$

So whatever we can do to find a $\lambda$ will give us our eigenvalues.
And that is why Arfken and other resources look for the conditions that make both of the recursion relations result in
a series solution that converges.
In the first case we need a value of lambda that will cancel the numerator $j(j+1)$ and in the second we need a value that
will cancel the $(j+1)(j+2)$.

\chapter{Variational Calculus}

Normed linear space - we want to define "closeness" via a norm and continuity to study functionals that look like
$$
J[y] = \int_{a}^{b} F(x, y, y^\prime) dx
$$

where the appropiate function space is the set of all functions with two continuous derivatives.
%%%%%%%%%%%%%%%%%%%%%%%%%%%%%%%%%%%%%%%%%%%%%%%%%%%%%%%%%%%%%%%%%%%%%%%%%%%%%%%
\subsubsection{A Propagator}

Describes the amplitude for a particle to propagate from one point to another.

With a source $j(\vect{x}, t)$, the Klein-Gordon equation becomes
$$
\left( \partial^2 + m^2 \right) \psi(\vect{x}, t) = j(\vect{x}, t)
$$

In momentum space the equation is
$$
\left( -\omega^2 + |\vect{k}|^2 + m^2 \right) \tilde{\psi}(\vect{k}, \omega) = \tilde{j}(\vect{\vect{k}, \omega})
$$

$$
\tilde{\psi}(\vect{k}, \omega) = \frac{ \tilde{j}(\vect{\vect{k}, \omega}) }{ \omega^2 - |\vect{k}|^2 - m^2 + i\epsilon }
$$

$i\epsilon$ Feynman's prescription.

The propagator $D_F(x-y)$ is the inverse transform of
$\frac{ 1 }{ \omega^2 - |\vect{k}|^2 - m^2 + i\epsilon }$

$$
D_F(x-y) = \int \frac{d^4k}{(2\pi)^4}
    \frac{ e^{-i(\vect{k}\cdot\vect{x} -\omega t)} }{ \omega^2 - |\vect{k}|^2 - m^2 + i\epsilon }
$$

\documentclass{article}
\usepackage{blindtext} % for table of contents and other niceseties.
\usepackage{amsmath}
\usepackage{amsfonts} % things such as \mathbb{R}, or \Z. See https://en.wikipedia.org/wiki/List_of_mathematical_symbols_by_subject
\usepackage{amstext} % for \text macro.
\usepackage{array} % for tables.
\usepackage{float} % here for H placement parameter
\usepackage{hyperref} % for hyperlinks.
\usepackage{listings} % for code.
\usepackage[makeroom]{cancel} % crossing out terms in math https://tex.stackexchange.com/questions/75525/how-to-write-crossed-out-math-in-latex
\usepackage[shortlabels]{enumitem} % enumerate with letters.


% Make greek letters bold too.
\usepackage{bm}
\newcommand{\vect}[1]{\boldsymbol{\mathbf{#1}}}

% Make paragrphs show in the table of contents.
\setcounter{secnumdepth}{4}
\setcounter{tocdepth}{4}

\title{Notes on Stats}

\numberwithin{equation}{section}

% There is a \pmod for the "mod" symbol, but not one for the "div" symbol.
% Thus we are making it up.
\makeatletter
\newcommand*{\bdiv}{%
  \nonscript\mskip-\medmuskip\mkern5mu%
  \mathbin{\operator@font div}\penalty900\mkern5mu%
  \nonscript\mskip-\medmuskip
}
\makeatother

\begin{document}

\maketitle
\tableofcontents

%%%%%%%%%%%%%%%%%%%%%%%%%%%%%%%%%%%%%%%%%%%%%%%%%%%%%%%%%%%%%%%%%%%%%%%%%%%%%%%%%%%%%%%%%%%%%%%
\section{Dynamic Properties of the Monte Carlo Method in Statistical Physics}

These are notes from the paper
"Dynamic Properties of the Monte Carlo Method in Statistical Physics"
by H. Muller-Krumbhaar and K Binder,
\href{https://link.springer.com/article/10.1007/BF01008440}{Journal of Statistical Physics, Vol 8, No. 1, 1973}.

References 5, 6, 27, and 28 mention processes for estimating statistical errors.

However, the processes mentioned in ref 28 disregard the knowledge that one can draw from
transition probabilities about the well-defined correlations between subsequent configurations.




%%%%%%%%%%%%%%%%%%%%%%%%%%%%%%%%%%%%%%%%%%%%%%%%%%%%%%%%%%%%%%%%%%%%%%%%%%%%%%%%%%%%%%%%%%%%%%%%%%%
\section{Central Limit Theorem}

A collection of random variables is independent and identically distributed (IID)
if each random variable has the same probability distribution as the others and
all are mutually independent.

Another way of thinking about an IID set of random variables is as a random sample.

If an arbitrarily large number of samples, each involving multiple observations (data points),
were separately used in order to compute one value of a statistic
(such as, for example, the sample mean or sample variance) for each sample,
then the sampling distribution is the probability distribution of the values that the statistic takes on.
\\

The central limit theorem establishes that, in many situations,
for independent and identically distributed random variables,
the sampling distribution of the standardized sample mean tends towards the standard normal distribution
even if the original variables themselves are not normally distributed.
\\

Normal distributions are often used to represent real-valued random variables
whose distributions are not known because of the central limit theorem.
One example of IID data could be measurement errors.




%%%%%%%%%%%%%%%%%%%%%%%%%%%%%%%%%%%%%%%%%%%%%%%%%%%%%%%%%%%%%%%%%%%%%%%%%%%%%%%%%%%%%%%%%%%%%%%%%%%
\section{Gaussians and the Error Function}

The error function was first formally introduced in
\href{https://books.google.com/books?id=8Po7AQAAMAAJ&pg=RA1-PA294#v=onepage&q&f=false}{On a class of definite integrals by J.W.L. Glaisher B.A. F.R.A.S. F.C.P.S.}.
(We found that link in the wikipedia page talking about the Error function
\href{https://en.wikipedia.org/wiki/Error_function}{Error function}, it was reference number 3.)

Some great references we recommend watching and working through are
\begin{enumerate}
\item \href{https://www.youtube.com/watch?v=jkytxdedxhU}{The impossible integral of e\^ (x\^ 2) \& the error function by blackpenredpen}
\item \href{https://www.youtube.com/watch?v=zorcLisjRUI}{Integral of e\^(x\^2) \& the Imaginary Error Function by blackpenredpen}
\end{enumerate}

$$
\erf(x) = \frac{2}{\sqrt{\pi}} \int_{0}^{x} e^{-t^2} dt
$$

This is the definite integral version of a Gaussian integral.
Since $\int_{-\infty}^{\infty} x^{-x^2} dx = \sqrt{\pi}$ and since the definite integral we use to define the
error function only goes from 0 to $x$ (only half the domain of the Gaussian),
that's why we normalize the $\erf(x)$ with a $\sqrt{\pi}/2$.
\\

To draw the horizontal asymptotes of $\erf(x)$, note that the domain is $\left(-\infty, \infty\right)$, so we need to take the
limits of $x$ as it tends to positive and negative infinity.

\begin{align*}
\lim_{x\rightarrow -\infty} \erf(x) &= \frac{2}{\sqrt{\pi}} \int_{0}^{-\infty} e^{-t^2} dt \\
&= - \frac{2}{\sqrt{\pi}} \int_{-\infty}^{0} e^{-t^2} dt \\
&= -1
\end{align*}

Similarly, $\lim_{x\rightarrow -\infty} \erf(x) = 1$.
\\

The derivative is also quite easy to get since we are taking the derivative of an integral and thus the fundamental theorem of calculus
is all we need,
$$
\erf^{\prime} (x) = \frac{2}{\sqrt{\pi}} e^{-x^2}
$$
From the derivative we can also tell that $\erf(x)$ has no critical numbers and as a matter of fact, the derivative is always non-negative.

And since we are talking about critical points, we can also go and take a look at the inflection points by computing the second
derivative,
$$
\erf^{\prime\prime}(x) = -2x \frac{2}{\sqrt{\pi}} e^{-x^2}
$$
From this expression we can tell that $x=0$ is an inflection point, and that for all $x<0$ $\erf(x)$ is concave down, and for $x>0$
$\erf(x)$ is concave up.
\\

Now, what about $\int e^{x^2} dx$, here we could do a variable transformation,
\begin{align*}
\int e^{x^2} dx &= \int e^{-\left(-x^2\right)} dx \\
& = \int e^{-(ix)^2} dx \rightarrow
\left[
    \begin{alignedat}{2}
        u  &= ix    \quad & x  &= \frac{1}{i}u = \frac{1}{i} \frac{i}{i} u = -iu \\
        du &= i dx  \quad & dx &= -i du
    \end{alignedat}\,
\right] \\
&= \int e^{-(u)^2} (-i) du \\
&= -i \int e^{-u^2} du \\
&= -i \frac{2}{\sqrt{\pi}} \erf(u) \\
&= -i \frac{2}{\sqrt{\pi}} \erf(ix) = \erfi(x)
\end{align*}

There is one more property to mention before we continue looking into variable transformations and that is that
$\erf(-z) = -\erf(z)$, the error function is an odd function.
This directly results from the fact that the integrand $e^{-t^2}$ is an even function
(the antiderivative of an even function which is zero at the origin is an odd function and vice versa).


\end{document}

\printbibliography[heading=bibintoc]


\end{document}